\documentclass[11pt]{article}
\usepackage{fontspec}
\usepackage{geometry}
\geometry{paperwidth=8.5in,paperheight=11in,left=1.25in,right=1.25in,top=1.5in,bottom=1.0in}
\setmainfont[
  BoldFont={KoPubWorldBatang Bold},
  ItalicFont={KoPubWorldBatang Medium},
  BoldItalicFont={KoPubWorldBatang Bold}
]{KoPubWorldBatang Medium}
\XeTeXlinebreaklocale "ko"
\XeTeXlinebreakskip = 0pt plus 1pt
\usepackage{amsmath,amssymb,amsfonts}
\usepackage{booktabs}
\usepackage{array}
\usepackage{hyperref}
\usepackage{xcolor}
\hypersetup{colorlinks=true,linkcolor=black,citecolor=black,urlcolor=black}
\usepackage{natbib}
\usepackage{caption}

\emergencystretch=2em

% Custom commands
\newcommand{\sgmvpi}{$(S, G, E, V, M, \pi)$}
\newcommand{\framework}{\textsc{Servo}}

% Theorem-like environments for the formal analysis
\usepackage{amsthm}
\newtheorem{proposition}{명제}
\newtheorem{definition}{정의}
\newtheorem{assumption}{가정}

% Title
\title{\textbf{AI 과학자 시스템의 형식화:\\
구성요소 프레임워크와 이론적 분석}}

\author{
  Jaesol Shin\\
  \small{WeDataLab}\\
  \texttt{ysys143@wedatalab.com}
}

\date{2026년 5월\\\small 2026년 7월 20일 개정}

\begin{document}

\maketitle

\begin{abstract}
가설 생성, 실험 실행, 지식 종합을 종단간(end-to-end)으로 자동화하는 시스템의 등장은 과학 수행 방식의 근본적 전환을 뜻한다. 빠른 경험적 진보에도 불구하고 최근 LLM 기반 \emph{AI 과학자} 시스템은 공유된 형식적 어휘 안에서 분석되는 경우가 드물어, 그 결과 서로 다른 아키텍처와 과학 분야 간의 원리적 비교가 어렵다. 본 논문은 \framework{}이라는 6요소 프레임워크 \sgmvpi{}를 제안한다. 이 프레임워크는 탐색 공간 \textbf{S}, 가설 생성기(\textbf{G}enerator), 실험 실행기(\textbf{E}xecutor), 검증기(\textbf{V}alidator), 메모리(\textbf{M}emory), 탐색 정책(\textbf{P}olicy, $\pi$)으로 구성되며, 부분 관측 마르코프 결정 과정(POMDP) 이론과 베이즈 실험 설계(BED)에 기반을 둔다. 우리는 이 프레임워크를 핵심 AI 과학자 시스템들과 일곱 개 과학 분야(수학, 물리학, 화학, 재료과학, 생물학, 사회과학, 컴퓨터과학) 사례에 적용하고, 그 시스템들의 출력 형식과 검증 계층을 다루는 보완적 연구 산출물 인프라 제안을 별도로 검토한다. 분석 결과 네 가지 지속적 미해결 과제가 드러난다. (1) 다층적 자동화 프록시에도 불구하고 인간 판단에 여전히 의존하는 참신성 평가의 비자동화 가능성, (2) 다루기 쉬운 근사가 이론적으로 가용함에도 BED 최적 탐색 정책의 부재, (3) 축적된 실험 경험이 향후 탐색을 개선하지 못하게 막는 의미 메모리 $M_s$의 미발달, (4) 계산적 대리 실험에서 물리적 검증으로 격상하는 원리적 기준의 부재, 곧 실험 충실도 간극이다. 우리는 검증기 $V$를 관측 우도와 구별하는 세 명제를 다루기 쉬운 식별 가정 하에서 형식화한다: 루프는 검증기와 무관하게 닫히고, 믿음 갱신에 잘못 결합된 검증기는 오설정된 표적으로의 수렴을 유발할 수 있으며, 모든 자동 채널에 부재하는 참신성 성분은 식별 불가능하다(명제~\ref{prop:suff}--\ref{prop:novelty}). 조사된 어떤 시스템도 이 가정들을 온전히 충족하지 않으므로, 이 명제들은 구분의 경험적 만연도를 입증하기보다 그 형식적 구조를 명료화한다. 교차 시스템 조사는 잠정적 구조화 주석으로 제시된다: 게이팅 검증기의 보정이 신뢰할 만한 폐쇄와 공기(共起)한다는 조직화 가설은 일부 결과 재코딩 가정 아래에서 방향이 일치하나 판정되지 않았으며, 경험적 발견으로 읽혀서는 안 된다.
\end{abstract}

\section{서론}

자동화된 과학적 발견---가설을 독립적으로 생성하고, 실험을 설계해 실행하며, 새로운 지식을 종합하는 시스템---에 대한 추구는 대규모 언어 모델(LLM)의 성숙과 함께 급격히 가속화되었다. Coscientist~\citep{boiko2023emergent} 같은 초기 시연은 LLM 기반 에이전트가 화학 합성을 자율적으로 계획하고 실행할 수 있음을 보여주었다. 이후 The AI Scientist~\citep{lu2024aiscientist}와 Agent Laboratory~\citep{schmidgall2025agentlab} 같은 시스템은 이 비전을 문헌 조사·가설 생성에서 코드 실행, 원고 작성, 자동 동료 심사에 이르는 연구 전체 주기로 확장했다. 가장 최근에는 Aletheia~\citep{google2026aletheia} 같은 시스템이 연구 수준의 수학 문제에 대응하기 시작했다. 생물의학 분야에서 Robin~\citep{ghareeb2026robin}은 문헌 탐색·가설 생성·RNA-seq 데이터 분석을 AI가 자율 수행하는 반자율 멀티에이전트 시스템을 시연했다(물리적 실험은 인간 연구자 담당). Co-Scientist~\citep{gottweis2026coscientist}는 Gemini 2.0 기반 협업 시스템으로 세 건의 생물의학 wet-lab 검증을 달성했다. 병행해서 Agent-Native Research Artifacts~\citep{liu2026lasthuman} 같은 제안은 보완적 문제를 다룬다. 즉, AI 과학자 시스템의 \emph{출력}이 어떤 형태여야 하며, 이를 어떻게 독립적으로 검증하느냐의 문제이다.

이런 경험적 추진력에도 불구하고, 이 분야에는 공유된 형식적 어휘가 결여되어 있다. 각 시스템은 고유한 용어로 자체 아키텍처를 도입하므로 체계적 비교가 어렵다. 실제 배포에 직결되는 질문들---\textit{어떤 요인이 시스템의 폐루프 작동 여부를 결정하는가? 자율적 참신성 평가가 일관되게 실패하는 이유는 무엇인가? 탐색 정책의 이론적 최적은 무엇인가?}---은 문헌에서 답을 얻지 못한 채 남아 있다.

우리는 다음 세 가지 기여로 이 공백을 메운다.

\begin{enumerate}
  \item AI 과학자 시스템을 6개 구성요소로 분해하고 이를 POMDP와 베이즈 실험 설계에 관한 이론적 문헌과 연결하는 형식적 프레임워크 \framework{}, \sgmvpi{}를 제안한다(\ref{sec:framework}절).
  \item \framework{}을 적용해 핵심 AI 과학자 시스템(\ref{sec:core}절), 보완적 연구 산출물 인프라 제안(\ref{sec:ara}절), 일곱 개 과학 분야의 도메인별 사례(\ref{sec:domains}절)를 분석한다.
  \item 모든 현재 시스템의 진보를 제약하는 네 가지 미해결 과제를 식별한다(\ref{sec:openproblems}절).
\end{enumerate}

\noindent\textbf{범위.} 본 서베이는 자동화된 과학적 \emph{발견}을 주된 목표로 삼는 시스템---가설-실험-검증 주기로 새로운 지식을 생성하는 시스템---을 다룬다. 인간이 수행하는 연구 파이프라인 안에서 AI를 특수한 도구로 적용하는 시스템(예: 구조 예측을 위한 AlphaFold~\citep{jumper2021alphafold})은 명시적으로 제외한다. 그런 시스템은 우리 프레임워크의 핵심인 가설-검증 루프가 없기 때문이다.

\section{배경}
\label{sec:background}

\subsection{과학적 발견의 POMDP 정식화}

우리는 AI 과학자 시스템을 \emph{부분 관측 마르코프 결정 과정}(Partially Observable Markov Decision Processes, POMDPs)~\citep{kaelbling1998pomdp}으로 모델링한다. 참 상태 $s \in \mathcal{S}$는 자연의 미지의 구조(예: 올바른 물리 법칙, 최적 분자 구성)에 해당한다. $s$에는 직접 접근할 수 없으므로, 시스템은 실험으로 얻은 관측값 $o \in \Omega$로 갱신되는 믿음 상태 $b(s)$를 유지한다.

형식적으로 AI 과학자 POMDP는 튜플 $(\mathcal{S}, \mathcal{A}, T, R, \Omega, \mathcal{O}, b_0, \gamma)$이다. 여기서 $\mathcal{A}$는 행동 공간(가설 생성, 실험 실행), $T: \mathcal{S} \times \mathcal{A} \rightarrow \Delta(\mathcal{S})$는 전이 함수, $R: \mathcal{S} \times \mathcal{A} \rightarrow \mathbb{R}$는 보상(reward; 발견의 과학적 가치), $\mathcal{O}: \mathcal{S} \times \mathcal{A} \rightarrow \Delta(\Omega)$는 잡음 있는 관측 과정, $b_0$는 초기 믿음(belief), $\gamma$는 할인 인자다. 믿음은 베이즈 규칙 $b'(s') \propto \mathcal{O}(o \mid s', a) \sum_{s} T(s' \mid s, a)\, b(s)$로 갱신되며, 정책 $\pi$는 접근 불가능한 상태 $s$가 아니라 믿음 $b$에 따라 행동한다.

이 정식화는 과학적 발견에 특유한 세 가지 측면에서 표준 POMDP와 다르다. (1) 보상 함수 $R$은 사전에 명시되지 않는다. 과학적 의의는 선험적으로 정의될 수 없기 때문이다. (2) 상태 공간 $\mathcal{S}$는 무한하고 개방적이다. (3) 전이 함수 $T$는 비정상적이다. 축적된 지식이 가치 있는 탐색의 경계를 변화시키기 때문이다. 따라서 우리는 POMDP를 개방형 과학의 완전히 명세된 모델이 아니라 조직화 형식으로 사용한다. \ref{sec:vformal}절의 형식 결과는 이 세 가지 일탈이 제한되는 다루기 쉬운 하위문제(가정~\ref{asm:tract})에서 성립하며, 일반적 개방형 경우는 미해결 문제로 남긴다.

\subsection{베이즈 실험 설계}

베이즈 실험 설계(Bayesian Experimental Design, BED)~\citep{chaloner1995bayesian,rainforth2024modern}는 탐색 정책 $\pi$의 원리적 기반을 제공한다. 최적 정책은 기대 정보 이득(Expected Information Gain, EIG)을 최대화하는 실험 $d^*$를 선택한다.

\begin{equation}
  d^* = \arg\max_{d \in \mathcal{D}}\; \mathrm{EIG}(d) = \arg\max_{d}\; \mathbb{E}\left[H(\theta) - H(\theta \mid y, d)\right]
  \label{eq:eid}
\end{equation}

여기서 $H(\theta)$는 매개변수 $\theta$에 대한 사전 엔트로피이고 $H(\theta \mid y, d)$는 실험 $d$ 하에서 결과 $y$를 관측한 후의 기대 사후 엔트로피다. 두 프레임워크는 경쟁적이라기보다 상호 보완적이다. POMDP가 결정 문제의 구조적 정식화를 맡고, BED는 정책 $\pi$의 규범적 기준을 맡는다.

최근 연구는 EIG가 변분 방법~\citep{foster2019variational}, 분할상환된 신경 정책~\citep{foster2021dad}, 강화학습~\citep{blau2022rl}로 근사된다는 것을 보였다. 이런 이론적 가용성에도 현재 어떤 AI 과학자 시스템도 EIG 최적 정책을 명시적으로 구현하지 않는다---우리는 이 공백을 핵심적인 미해결 과제로 식별한다. 다만 이 진술은 좁혀 읽어야 한다. 도메인 특화 시스템은 이미 능동학습 근사를 사용한다---Robot Scientist의 실험 선택~\citep{king2004robot}과 GNoME의 DFT-in-the-loop 능동학습~\citep{merchant2023gnome}이 그 예다. 우리가 지적하는 것은 그보다 좁은 간극이다: 최근의 개방형 LLM 기반 시스템은 명시적 사후분포도 진술된 EIG 목적도 유지하지 않는다.

\subsection{검증기 완결성: 형식적 분석}
\label{sec:vformal}

위의 정식화는 검증기 $V$를 관측 우도와 구별해 취급할 때 무엇이 따라 나오는지를 정밀하게 진술하게 해준다. 우리는 $V$를 보상 채널로 취급하고, POMDP 가정이 성립하는 다루기 쉬운 하위문제에서 루프 폐쇄를 분석한다. 개방형 영역은 미해결 문제로 남긴다.

\begin{assumption}[다루기 쉬운 하위문제]
\label{asm:tract}
(i) 탐색 공간 $S$가 유계이고 모수 공간 $\Theta$가 콤팩트(compact)하다; (ii) 전이 $T$가 정상적이다; (iii) 사전 분포 $p(\theta)$의 지지집합이 참값 $\theta^\star$를 포함한다; (iv) 관측 모델 $p(o \mid \theta, d)$가 올바르게 명세되어 있으며 $\theta$를 식별한다($\theta \neq \theta'$인 두 값에 대해 이를 분리하는 도달 가능한 설계가 존재한다); (v) 실현된 설계 열이 사후 일치성에 대해 \emph{충분히 정보적}이다---무한히 자주, 살아남은 임의의 두 가설 사이의 쿨백--라이블러 분리가 고정된 $\delta>0$ 이상인 도달 가능한 설계를 선택한다. 우리는 이를 적응적 EIG 항의 발산하는 합과 동일시하지 않는다: 증분 EIG는 유한 상태 설정에서 초기 엔트로피로 상계되는 기대 엔트로피 감소량이므로, 사후 분포가 집중하더라도 누적 EIG는 유계일 수 있다. 여기서의 조건은 누적 EIG가 아니라 가설 간 KL 분리에 대한 것이다.
\end{assumption}

\begin{definition}[보상 채널로서의 검증기]
\label{def:channel}
관측 $o \sim \mathcal{O}(\cdot \mid s,a)$에 대해, 검증기는 참 보상 $R(s,a)$의 추정치 $\hat{r} = V(o)$를 유도한다. $\mathbb{E}[V(o) \mid s,a] = R(s,a)$이면 $V$를 \emph{보정됨(calibrated)}, 어떤 $(s,a)$에서 계통 오차 $\delta(s,a) = \mathbb{E}[V(o) \mid s,a] - R(s,a)$가 0이 아니면 \emph{편향됨(biased)}이라 한다. $V$는 명시적으로 결합되지 않는 한 갱신 가능도 $p(o \mid \theta, d)$와 구별된다.
\end{definition}

\begin{definition}[루프 폐쇄]
\label{def:closure}
정책 $\pi$ 하에서 $t \to \infty$일 때 믿음 $b_t(\theta)$가 참값 $\theta^\star$에 집중되면 발견 루프가 $\theta$에서 \emph{닫힌다}.
\end{definition}

\begin{proposition}[폐쇄를 위한 충분조건]
\label{prop:suff}
가정~\ref{asm:tract} 하에서 사후 분포는 $\theta^\star$에 집중하며(루프가 닫힘), 일단계 EIG 규칙(\ref{eq:eid}식)은 정보 목적에 대해 근시안적으로 최적이다---둘 다 \emph{검증기와 무관하게} 성립한다. $V$는 보상 채널인 반면 EIG는 사전 분포와 가능도만의 범함수이기 때문이다.
\end{proposition}

\noindent\emph{증명 개요.} 올바르게 명세되고 식별하는 가능도가 콤팩트한 $\Theta$ 위에서 \emph{발산하는} 누적 정보를 가지는 것(가정~\ref{asm:tract}(i),(iv),(v))은 베이즈 사후 일치성의 표준 충분조건이며~\citep{ghosal2017fundamentals}, $b_t \to \delta_{\theta^\star}$를 준다. 보정은 일치성에도 EIG 최적성에도 요구되지 않는다---둘 다 사전 분포와 가능도에만 의존한다. 보정은 별개의 보상 채널에 관한 것이며, 이는 아래에서 정밀화한다. 균일 정보 조항 (v)는 본질적이다: 모든 $t$에서 $\mathrm{EIG}(d_t)>0$이면서 $\sum_t \mathrm{EIG}(d_t)<\infty$인 열은 식별 가능성과 보정이 모두 성립하더라도 사후 분포를 점질량에 이르게 하지 못한다. $\square$

\begin{proposition}[오설정된 결합 하의 표류]
\label{prop:drift}
검증기가 믿음 갱신에 진입한다고 하자. 즉 실효 갱신이 결합 계수 $\beta>0$에 대해 $q_\theta(o) \propto p(o \mid \theta, d)\,\exp(\beta V(o))$라 하자. $V$가 결과에 $\theta$-의존적인(비아핀) 방식으로 의존하면 사후 분포는 쿨백--라이블러(KL) 사영 $\theta^\dagger$에 집중하며~\citep{berk1966}, 이는 일반적으로 $\theta^\star$와 다르다: 루프는 계산적으로 닫히지만 다른 표적 위에서 닫힐 수 있다. $V$의 $\theta$-무관 가법 상수는 로그 분배함수에서 상쇄되어 $\theta^\star$를 바꾸지 않으며(척도 변화는 대신 결합 계수 $\beta$를 재조정한다), $V$가 순수한 정책 채널일 때($\beta=0$)는 어떤 검증기 변환도 표적을 옮기지 못한다---기껏해야 설계 선택을 악화시키고 정보 획득을 늦출 뿐이다(가정~\ref{asm:tract}(v)와 상호작용한다).
\end{proposition}

\noindent\emph{증명 개요.} 결합된 오설정 갱신 가능도는 사후 분포를 그 KL 최소점에 놓는다~\citep{berk1966}. 로그 분배함수 $\log\!\int p(o\mid\theta,d)\,e^{\beta V(o)}\,do$는 그 $\theta$ 의존성이 바뀔 때에만 최소점을 이동시키는데, $\theta$-무관 가법 상수는 이를 바꾸지 않는다. $\beta=0$이면 갱신 가능도가 불변이므로 $\theta^\star$는 $V$에 대해 불변이다. $\square$

\begin{proposition}[참신성의 비식별성]
\label{prop:novelty}
식별되는 성분들과 사전적으로 독립인 참신성/유의성 성분을 $\theta_{\mathrm{nov}}$라 하자. 모든 자동 검증기의 채널 가능도가 $\theta_{\mathrm{nov}}$에 대해 독립이면, 어떤 자동 검증기도 $\theta_{\mathrm{nov}}$의 일치 추정량을 갖지 못하며, 참신성에 대한 루프 폐쇄는 $\theta_{\mathrm{nov}}$에 관해 정보를 주는 추가 채널을 요구한다---현재로서는 $V_{\mathrm{human}}$만이 이를 공급한다.
\end{proposition}

\noindent\emph{증명 개요.} 가능도가 $\theta_{\mathrm{nov}}$에 의존하지 않고 사전 분포가 이를 식별되는 성분들과 결합시키지 않으면, $\theta_{\mathrm{nov}}$의 사후 주변 분포는 모든 데이터에 대해 사전 분포와 같다. 따라서 $\theta_{\mathrm{nov}}$는 식별되지 않으며 일치 추정량이 존재하지 않는다. $\square$

이 형식 결과들은 더 좁은 구분을 뒷받침한다: 다루기 쉬운 하위문제에서의 사후 폐쇄는 올바르게 명세되고 식별하며 충분히 정보적인 관측 모델에 의존하는 반면(명제~\ref{prop:suff}), 검증기 보정은 정책이 사용하는 보상 또는 수용 신호를 지배한다---사후 일치성을 지배하지 않으며, 어떤 설계가 EIG를 최대화하는지도 지배하지 않는다(정보 최적 설계와 보상 최적 설계는 검증기 보상이 EIG 순서를 보존할 때에만 일치하며, 완벽히 보정된 $V$에 대해서도 갈라질 수 있다). 명제~\ref{prop:drift}는 표류가 단순한 순위 변동이 아니라 잘못 \emph{결합된} 검증기를 요구함을 보이고, 명제~\ref{prop:novelty}는 어떤 자동 검증기도 해결할 수 없는 유일한 성분을 분리한다. 이 결과들만으로는 경험적 $V$ 완결성 순서가 모든 신뢰할 만한 자율 발견에 필요하다는 것이 증명되지 않는다; 그 주장은 독립적으로 코딩된 결과를 요구하는 설계 가설로 남는다. 따라서 조사된 시스템들은 관측, 보상 검증, 참신성 판단을 분리하기 위한 동기부여 사례로 읽혀야 하며, $V$ 완결성이 폐쇄의 일차적 결정 요인이라는 형식적 논증으로 읽혀서는 안 된다.

\section{관련 연구}
\label{sec:related}

\textbf{AI 과학자 서베이 및 시스템 논문.} AI 과학자 시스템 문헌은 빠르게 증가했지만 공유된 어휘가 부족하다. \citet{lu2024aiscientist}은 최초의 폐루프 종단간 시스템을 제안하며 자체 구성요소에 대한 정성적 특성화를 제공하지만, 시스템 간 비교를 위한 형식적 프레임워크를 제시하지는 않는다. \citet{schmidgall2025agentlab}은 코파일럿 설계에 초점을 두고 인간 기준선 대비 비용 절감을 평가한다. \citet{boiko2023emergent}은 도구로 강화된 자율 합성을 시연하지만 발견 루프를 모델링하지는 않는다. \citet{si2024novelideas}은 LLM 아이디어 생성 품질을 대상으로 한 대규모 인간 연구를 수행해 참신성 평가의 기준선을 수립했다. \citet{gu2024scimuse}는 58M 논문 지식 그래프 기반 개인화 아이디어 생성 시스템(SciMuse)을 제안하고 110명의 Max Planck 연구그룹 리더가 4,451개 제안을 평가한 최대 규모의 인간 검증을 수행했다; 특히 PageRank와 피인용 수가 높은 개념일수록 흥미도가 낮다는 음의 상관을 실증해, impact 지표 기반 탐색이 흥미로운 아이디어 발굴과 오히려 역행함을 경험적으로 보였다. \citet{artiles2026alienspace}은 이 한계의 원인을 인지적 가용성(cognitive availability) 편향으로 정량 진단하고, coherence 모델과 availability 모델을 분리 훈련해 인간 연구자가 인지적으로 접근하기 어려운 연구 방향을 프런티어 LLM 대비 3.5--7배 넓은 탐색 공간에서 샘플링하는 프레임워크를 제안한다. \citet{lee2025spacer}는 이와 유사하게 \textit{의도적 탈맥락화}(deliberate decontextualization)를 제안해, 180,000편 생물학 논문을 키워드 집합으로 분해하고 맥락 제거 상태에서 새로운 연결을 탐색하며, 임베딩 공간에서 SOTA LLM보다 선도 논문에 더 가까운 아이디어를 생성한다. \citet{chen2026gap}는 평가의 단위를 개별 아이디어에서 \emph{분포}로 옮겨, 고품질 논문의 핵심 아이디어를 촉발했을 선행연구 소집합을 역설계한 뒤 2축 research-taste taxonomy로 인간 참조 분포와 LLM 생성 분포의 divergence를 계측한다; 여러 LLM에 걸쳐 LLM 아이디어가 bridge형 기회와 synthesis 방법에 체계적으로 집중되어 인간 연구 취향보다 좁고 편향되어 있음을 실증해, 참신성 자동화의 한계(명제~\ref{prop:novelty})에 분포 수준의 경험적 근거를 더한다. \citet{baek2024researchagent}은 문헌 조건부 반복 아이디에이션을 제안한다. 이들 중 어느 것도 탐색·실행·검증·메모리를 통합하는 프레임워크를 제공하지는 않는다. 우리 프레임워크와 가장 가까운 선행 연구는 \citet{lu2026aiscientist}의 정성적 논의로, 개방형 아카이브와 재현 노드 개념을 도입하지만 이를 일반 시스템의 구성요소로 형식화하지는 않는다. \citet{ghareeb2026robin}은 생물의학 도메인에서 $V_\text{empirical}$ 분석 루프(RNA-seq DGE 분석)를 AI가 자율 처리하는 최초의 멀티에이전트 시스템을 시연하며(물리적 실험 실행은 인간 담당), \citet{gottweis2026coscientist}는 Elo 토너먼트 기반 가설 진화와 세 건의 독립적 wet-lab 검증을 달성한다. 두 논문 모두 \framework{}의 $V_\text{empirical}$ 구성요소가 생물의학 과제에서 실질적으로 작동함을 처음으로 보인다. \citet{aygun2026era}는 LLM과 트리 탐색을 결합해 전문가 수준의 경험적 소프트웨어를 자동 생성하며, 이는 \framework{}의 $E$ 충실도 스펙트럼에서 계산 실행을 전문화한 시스템이다. \citet{liu2026autoresearchclaw}은 자기치유 실행기, 크로스-프로젝트 실패 메모리, 다중 에이전트 토론, 7단계 인간 개입 모드를 통합해 ARC-Bench에서 AI Scientist v2 대비 54.7\% 향상을 달성하며, \framework{}의 $H$ 매개변수 설계 공간에 최초의 경험적 증거를 제공한다. \citet{meng2026scientistone}은 Chain-of-Evidence(CoE) 프레임워크를 통해 모든 주장에 증거 출처를 강제 연결하고, 4중 CoE Audit으로 환각 인용 0/337을 달성하며, \framework{}의 $V$ 계층---$V_\text{syntax}$와 $V_\text{semantic}$ 동시 자동화---에 대한 가장 완결된 구현을 제시한다. \citet{gao2026autoscientists}은 분산 자기조직화 멀티에이전트 시스템 AutoScientists를 제안해 6개 생의학 도메인에서 병렬 탐색과 공유 실패 메모리로 단일 에이전트 대비 7개 대 0개 개선을 달성하며, \framework{}의 $M_e \to M_s$ 변환과 분산 $\pi$ 설계에 실증적 근거를 제공한다.

\textbf{유사 형식화 시도.} 본 연구와 가장 가까운 입장을 취하는 선행 연구들이 존재하나, 각각 \framework{}가 다루는 범위의 일부만을 포괄한다. \citet{tie2025survey}는 문헌 조사·아이디어 생성·실험 준비·실험 실행·글쓰기·논문 생성의 6단계 파이프라인으로 수십 편의 AI 과학자 시스템을 분류하지만, POMDP나 BED에 기반한 수학적 형식화나 $V$/$M$ 계층 분해를 제공하지 않는다. \citet{tie2026autoresearch}는 동일 연구 그룹이 23인 저자로 확장한 후속 서베이로, L0--L4 자율성 스펙트럼과 참신성·타당성·영향력·신뢰성·출처의 5가지 과학적 신뢰성 차원을 제안하며 인간 주도 보조(L1--L2)를 ``Vibe Research''로 명명하지만, 마찬가지로 \framework{}가 제공하는 수학적 $V$/$M$ 계층 분해나 POMDP-BED 기반 형식화는 다루지 않는다. \citet{exhyte2025}는 탐색-가설-검증(EXHYTE) 주기를 중심으로 83편을 정리하는 프로세스 중심 서베이로, 직관적 단계 분류에 머무른다. \citet{hypoagents2025}는 베이즈 정리와 정보 엔트로피를 결합해 가설 집합을 반복 최적화하는 HypoAgents를 제안하며 $G$ 구성요소를 부분적으로 다루나, 전체 $E$-$V$-$M$-$\pi$ 사이클을 포괄하지 않는다. \citet{pievo2026}은 가우시안 프로세스 기반 정보 지향 표집(Information-Directed Sampling)으로 확장되는 원리 공간에서 베이즈 최적화를 수행하며 SOTA 대비 29.7\%--31.1\% 향상을 보이지만, 이 역시 $\pi$(가설 탐색 정책) 단계에 한정된 알고리즘 기여다. 본 논문은 이 네 연구가 각기 다루는 단계들---서베이의 분류 체계, $G$의 베이즈 최적화, $\pi$의 정보 이득 표집---을 \framework{}라는 단일 형식적 프레임워크 안에서 통합하고, 여기에 $V$ 4층 분해, $M$ 3층 분해, $H$ 매개변수, 7개 도메인 병목 진단을 추가한다.

\textbf{이론적 기반: POMDP와 BED.} 부분 관측 하 의사결정의 POMDP 정식화는 고전적이다~\citep{kaelbling1998pomdp}. 베이즈 실험 설계는 최적 실험 선택 기준을 형식화한다~\citep{chaloner1995bayesian}. 최근 연구는 변분 추론~\citep{foster2019variational}, 분할상환된 신경 정책~\citep{foster2021dad}, 강화학습~\citep{blau2022rl}로 다루기 쉬운 근사를 제공한다. 우리는 이 두 흐름의 연구를 현대 LLM 에이전트 환경에 대해 연결한다. POMDP가 구조적 정식화를 맡는 한편 BED는 규범적 정책 기준을 맡는다. 베이즈 최적 실험 선택은 자동 과학에서 오랜 역사를 가지며---가장 직접적으로는 정보 이득이 최대인 실험을 선택하는 능동학습 정책~\citep{king2004robot}을 갖춘 Robot Scientist~\citep{sparkes2010robot}에서---우리는 이 아이디어 자체의 선취권을 주장하기보다 이 전통 위에 구축한다.

\textbf{자율 과학 프레임워크.} 폐루프 자율 발견은 LLM 시대에 선행한다. Robot Scientist~\citep{sparkes2010robot}는 효모 기능유전체학에서 가설 생성, 로봇 실험, 통계적 검증을 통합하여, LLM 기반 시스템보다 10년 이상 앞서 \framework{}의 여섯 구성요소 전체---능동학습 정책~\citep{king2004robot} 포함---를 물리적 웻랩 형태로 사례화했다. 보다 최근의 통합 프레임워크로는 아이디어 생성·방법론 구성·실험 실행·결과 피드백을 12개 과학 과제에 걸쳐 잇는 폐루프 다중 에이전트 시스템 NovelSeek~\citep{zhang2025novelseek}, 그리고 생물의학 wet-lab 검증을 갖춘 다중 에이전트 가설 생성 시스템 Co-Scientist~\citep{gottweis2026coscientist}가 있다. 따라서 우리는 \framework{}을 자동 발견을 위한 최초의 통합 프레임워크가 아니라, 검증기 완결성·의미 메모리·POMDP/BED 관계를 최근 LLM 에이전트 하위 분야 안에서 명시화하는 형식적 어휘로 위치시킨다.

\textbf{도메인 특화 AI 과학 시스템.} 일반적 프레임워크를 제안하지 않으면서 특정 도메인의 과학 문제에 AI를 적용하는 상당한 선행 연구가 존재한다. 수학에서는 FunSearch~\citep{romera2023funsearch}와 DeepSeek-Prover~\citep{xin2024deepseek}, 물리학에서는 AI Feynman~\citep{udrescu2020afeynman}, 화학에서는 ChemCrow~\citep{bran2023chemcrow}와 ChemReasoner~\citep{sprueill2024chemreasoner}, 생물학에서는 BioPlanner~\citep{odonoghue2023bioplanner}와 실험 검증까지 수행한 멀티에이전트 Virtual Lab~\citep{swanson2025virtuallab}(LLM PI가 이끄는 역할 분담 에이전트 팀이 ESM·AlphaFold-Multimer·Rosetta 파이프라인으로 92개 나노바디를 설계·wet-lab 검증), 컴퓨터과학에서는 AutoML-Zero~\citep{real2020automlzero}가 있다. 우리는 이런 시스템들을 독립적 기여라기보다 \framework{} 구성요소의 사례로 다루며, 그렇게 함으로써 도메인 간의 원리적 비교를 가능케 한다.

\textbf{과학의 형식적 모델.} \citet{taniguchi2024cpc}는 과학 활동을 집단 예측 코딩(Collective Predictive Coding as Model of Science, CPC-MS)으로 재정식화한다. 다양한 사전 분포 $\theta_k$를 가진 에이전트들이 실험과 동료 심사로 공유 표현 $w$를 갱신하는 분산 베이즈 추론으로 과학을 모델링한다. 이 프레임워크는 \framework{}의 구성요소에 직접 대응한다: 가설 생성이 내적 표현 추론 $q(z^k \mid w, o^k)$, 실험이 행동 조건부 관측 $p(o \mid z, a)$, 동료 심사가 집단 합의 $q(w \mid \{z^k\})$에 해당한다. 핵심 차이는 범위다: CPC-MS는 \emph{집단적·사회적} 모델인 반면 \framework{}는 \emph{단일 시스템·모듈 분해}다. 특히 CPC-MS는 Lu et al.의 AI Scientist를 직접 분석하며, 이 시스템이 단일 연구자의 루프만 자동화할 뿐 협업적 동료 심사와 합의 형성이라는 사회적 계층은 다루지 않는다고 주장한다---\ref{sec:discussion}절에서 다루는 집단적 확장의 동기가 된다.

\textbf{연구 산출물 인프라.} \citet{liu2026lasthuman}은 AI가 생성한 연구의 구조적 출력 형식이자 검증 프로토콜로 Agent-Native Research Artifacts(ARA)를 제안한다. 우리는 이 기여를 종단간 AI 과학자 시스템과 별도로 위치시킨다. 발견 루프가 아닌 연구 산출물 계층을 다루기 때문이다. \framework{}의 $V$ 및 $M_s$ 구성요소에 대한 함의는 \ref{sec:ara}절과 \ref{sec:openproblems}절에서 논의한다.

\section{\framework{} 프레임워크}
\label{sec:framework}

우리는 AI 과학자 시스템을 다음 구성요소로 이루어진 튜플 \sgmvpi{}로 정의한다.

\begin{equation}
  \text{\framework{}} = (S,\; G,\; E,\; V,\; M,\; \pi)
\end{equation}

\begin{description}
  \item[$S$: 탐색 공간] 탐색 대상이 되는 후보 가설, 프로그램, 분자, 실험 조건, 그 밖의 객체 집합. $S$는 정적이거나 동적으로 확장(개방형)될 수 있다.

  \item[$G: M \times S \rightarrow h$: 가설 생성기] 축적된 메모리 $M$과 탐색 공간의 현재 상태에 조건부로 새로운 후보 $h \in S$를 제안하는 함수. $G$는 LLM 생성, 진화적 변이, 기호적 탐색을 포함할 수 있다.

  \item[$E: h \times P \rightarrow o$: 실험 실행기] 가설 $h$를 프로토콜 $P$ 하에서 실행해 관측값 $o$를 산출하는 함수. $E$의 충실도는 빠른 계산 시뮬레이션(저비용, 잠재적 대리자 오차)에서 물리적 웻랩 실행(고비용, 진실값)에 이른다.

  \item[$V: o \rightarrow r$: 검증기] 관측값 $o$에 보상 신호 $r$을 부여하는 함수. 우리는 $V$의 네 계층을 구분한다.
    \begin{align*}
      V_\text{syntax}   &: o \rightarrow \{\text{통과}, \text{실패}\} & \text{(결정론적, 빠름)}\\
      V_\text{semantic} &: o \rightarrow r \in [0,1]                  & \text{(LLM 기반, 확률적)}\\
      V_\text{empirical}&: o \rightarrow r                            & \text{(재현 기반)}\\
      V_\text{human}    &: o \rightarrow \{\text{참신함}, \text{유의함}\} & \text{(자동화 불가)}
    \end{align*}

  \item[$M = (M_e, M_s, M_p)$: 메모리] 실험 이력의 구조적 축적. $M_e$(에피소드 메모리)는 개별 실험 기록 $(h_i, o_i, r_i)$를 저장하고, $M_s$(의미 메모리)는 정제된 지식 패턴을, $M_p$(절차 메모리)는 학습된 실험 프로토콜을 저장한다. $M$의 구조가 $G$와 $\pi$의 품질을 직접 결정한다.

  \item[$\pi: b(S) \times M \rightarrow h$: 탐색 정책] 다음 가설을 선택하는 의사결정 규칙. 믿음 상태 $b(S)$는 매 실험 이후 갱신된다. BED 해석에서 $\pi^* = \arg\max_d \mathrm{EIG}(d)$이다(\ref{eq:eid}식). 실제로 현재 시스템들은 반응적(ReAct)에서 탐욕적, 트리 탐색에 이르는 정책을 사용한다.
\end{description}

\noindent 운영 루프는 다음과 같다.
\begin{equation}
  \pi(b, M) \;\rightarrow\; G(M_s, S) \;\rightarrow\; E(h, P) \;\rightarrow\; V(o) \;\rightarrow\; M \mathrel{+}= (h, o, r) \;\rightarrow\; b(S) \text{ 갱신} \;\rightarrow\; \text{반복}
\end{equation}

인간 개입 정도를 나타내는 보조 매개변수로 $H \in [0,1]$을 도입한다. 완전 자율 시스템에서는 $H=0$, 인간이 $G$나 $\pi$의 역할을 떠맡는 시스템에서는 $H=1$이다.

\section{핵심 AI 과학자 시스템 분석}
\label{sec:core}

종단간 자동화 연구 파이프라인을 제안하거나 구현한 4개의 핵심 시스템(\ref{sec:coscientist}--\ref{sec:agentlab}절)에 \framework{}을 적용하고, 그 시스템들의 출력 형식과 검증 계층을 다루는 하나의 연구 산출물 인프라 제안(\ref{sec:ara}절)을 검토한다. 각 종단간 시스템에 대해 6개 구성요소를 특성화하고 세 가지 품질 차원---(A) 참신성 주장 메커니즘, (B) 폐루프 완결성, (C) 과학적 품질 보증---을 평가한다.

\subsection{Coscientist}
\label{sec:coscientist}

\citet{boiko2023emergent}은 GPT-4가 도구로 강화된 ReAct 루프를 통해 화학 합성을 자율적으로 계획하고 물리적으로 실행할 수 있음을 시연한다.

\begin{table}[h]
\centering
\small
\begin{tabular}{@{}lp{9cm}@{}}
\toprule
구성요소 & 사례 \\
\midrule
$S$ & 알려진 화학 반응 공간(재현, 발견 아님) \\
$G$ & 인간 제공 과제 $\rightarrow$ GPT-4 Planner 단계; 자율적 아이디에이션 없음 \\
$E$ & 도구 호출(웹, 문서, 코드) + RoboRXN 물리 합성; 분석적 특성화는 인간에 위임 \\
$V$ & $V_\text{human}$ + 제한적 $V_\text{empirical}$: 저자 검사와 GC-MS; 반응 최적화 과제는 측정 결과를 과제 수준 피드백으로 활용 \\
$M$ & 세션 대화 이력; 세션 간 지속성 없음 \\
$\pi$ & 탐욕적 ReAct: 현재 문맥에 조건부인 다음 도구 행동; 분기 없음 \\
\bottomrule
\end{tabular}
\caption{Coscientist의 \sgmvpi{} 사상.}
\end{table}

\noindent\textbf{(A) 참신성.} 참신성 측정 메커니즘이 없다. 모든 합성 표적은 알려진 화합물이다. 시스템의 참신성 주장은 과학적 발견이 아니라 아키텍처적 능력에 관한 것이다.

\noindent\textbf{(B) 루프.} 루프는 \emph{부분적}이다. Coscientist는 제한된 과제 수준 피드백을 보인다: 반응 최적화 실험에서 측정 결과가 이후 선택에 반영되며, 정규화된 성능이 반복에 걸쳐 향상된다~\citep{boiko2023emergent}. 닫히지 않는 것은 참신성 수준의 발견 루프다---새로운 과학적 지식을 자율적으로 인증하는 가설 $\rightarrow$ 실험 $\rightarrow$ 새 가설 주기는 없으며, 참신성과 유의성의 검증은 인간이 매개한다.

\noindent\textbf{(C) 품질.} 단일 GC-MS 측정이며 재현, 통계적 검정, 공개 데이터 배포가 없다(안전상의 이유로 보류). 저자들은 다중 실행 분산 분석을 향후 과제로 남긴다.

\subsection{The AI Scientist (Sakana AI)}
\label{sec:aisci-sakana}

\citet{lu2024aiscientist}은 아이디에이션·실험·원고 작성·자동 동료 심사를 포함해 전체 연구 루프를 닫는 최초의 시스템을 제안한다.

\begin{table}[h]
\centering
\small
\begin{tabular}{@{}lp{9cm}@{}}
\toprule
구성요소 & 사례 \\
\midrule
$S$ & ML 연구 아이디어 공간; 시드 코드 템플릿이 암묵적으로 한정 \\
$G$ & 아카이브 $M_e$에 조건부인 LLM 아이디에이션; Semantic Scholar API(10라운드) 유사도 필터 \\
$E$ & Aider 코드 에이전트 $\rightarrow$ Python 실행; 최대 4회 재시도, 7200초 타임아웃 \\
$V$ & $V_\text{semantic}$: GPT-4o 리뷰어(NeurIPS 가이드라인, 5-앙상블); BalAcc 0.65(ICLR 2022 accept/reject 라벨 예측)---인간 상호일치 0.66은 다른 표본에서 측정된 별개의 량이다 \\
$M$ & $M_e$: 제목, 점수, 결과를 담은 아이디어 아카이브; $M_s$나 $M_p$ 없음 \\
$\pi$ & 탐욕적: 아카이브와 비유사한 가설 생성; 장기 계획 없음 \\
\bottomrule
\end{tabular}
\caption{The AI Scientist (Sakana 2024)의 \sgmvpi{} 사상.}
\end{table}

\noindent\textbf{(A) 참신성.} LLM이 각 아이디어의 참신성을 1--10 척도로 자체 점수화한다. 저자들은 체계적 과대평가 편향과 ``LLM 판단은 종종 편향을 가질 수 있다''~\citep{lu2024aiscientist}는 점을 인정한다. 중요하게도 저자들은 이렇게 진술한다. \textit{``우리는 이 버전의 The AI Scientist의 과학적 내용을 액면 그대로 받아들일 것을 권장하지 않는다.''}

\noindent\textbf{(B) 루프.} 진정으로 닫힌 최초의 루프다. 심사 점수가 $M_e$를 거쳐 다음 아이디에이션 주기로 피드백된다. 그러나 $\pi$는 탐욕적이며 아카이브 성장과 함께 개선되지 않는다.

\noindent\textbf{(C) 품질.} 통계적 엄밀성이 약하다. 시드와 신뢰구간이 강제되지 않는다. 초기 시스템 버전은 ablation 표를 환각으로 생성했다. Python 프로세스를 무한히 생성하는 서브프로세스와 자체 타임아웃을 연장하려 시도한 에이전트를 포함한 안전 위반이 관찰되었다.

\subsection{The AI Scientist (Nature 2026)}
\label{sec:aisci-nature}

\citet{lu2026aiscientist}은 Sakana 시스템을 에이전트 기반 트리 탐색과 개방성 아카이브로 확장하여, 동료 심사 워크숍(ICLR 2025 ICBINB, 점수 6.33)에서 수락된 최초의 AI 저작 원고를 달성한다(사전 합의에 따라 출판 전 철회).\footnote{The AI Scientist의 (2024)판과 (2026)판은 각각 이 시스템의 템플릿 기반 모드와 에이전트 기반 트리 탐색 모드(``AI Scientist-v2''~\citep{yamada2025aiscientistv2})이며, 후자가 동료 심사를 거친 \textit{Nature} 판이다~\citep{lu2026aiscientist}.}

\begin{table}[h]
\centering
\small
\begin{tabular}{@{}lp{9cm}@{}}
\toprule
구성요소 & 사례 \\
\midrule
$S$ & 개방성 아카이브(동적 $S$); Semantic Scholar(20라운드)에 대한 아이디어 참신성 검증 \\
$G$ & $S_\text{archive}$와 $M_e$에 조건부인 LLM; Semantic Scholar 유사도 필터 \\
$E$ & 4단계 에이전트 트리 탐색; 그림에 대한 VLM 비평; 4회 재시도의 Aider \\
$V$ & $V_\text{empirical}$: 재현 노드(평균 $\pm$ 표준편차); $V_\text{semantic}$: 자동 리뷰어(BalAcc 0.69); $V_\text{human}$: ICLR 리뷰어 \\
$M$ & $M_e$: 트리 노드; $M_s$: 개방성 아카이브; 재현 통계 \\
$\pi$ & 최선 우선 트리 탐색: GPT-4o가 메트릭, 학습 동역학, 그림 품질로 노드 점수화 \\
\bottomrule
\end{tabular}
\caption{The AI Scientist (Nature 2026)의 \sgmvpi{} 사상.}
\end{table}

\noindent\textbf{(A) 참신성.} 3계층 검증: 자기 점수, API 필터, 실제 동료 심사. 자동 리뷰어 성능: 균형 정확도 0.69(인간 0.66 대비), F1 0.62(인간 0.49 대비), 사전 컷오프($p=0.319$) 및 사후 컷오프($p=0.921$) 데이터 모두에서 인간 정확도와 유의한 차이 없음. 한계가 남아 있다: 환각 인용과 사전학습 데이터가 참신성 판단을 오염시키는 본질적 위험.

\noindent\textbf{(B) 루프.} 현재까지 가장 완전한 폐루프다. 트리 탐색이 보상 유도 $\pi$를 만들고, VLM 비평이 $E$ 내 피드백을 만들며, 재현 노드가 통계적 검증을 강제한다.

\noindent\textbf{(C) 품질.} 부트스트랩 신뢰구간(5,000회 복제)과 z-검정이 보고된다. 환각 인용은 문서화된 실패 양상으로 남아 있다. 워크숍 수락률(70\%)이 본 학회 수락률(32\%)과 다르므로, 동료 심사 검증의 의미는 제한적이다.

\subsection{Agent Laboratory}
\label{sec:agentlab}

\citet{schmidgall2025agentlab}은 인간 연구자가 연구 방향을 정하고 LLM 에이전트가 연구 파이프라인(문헌 검토, 실험, 원고 작성)을 실행하는 코파일럿 설계를 채택한다.

\begin{table}[h]
\centering
\small
\begin{tabular}{@{}lp{9cm}@{}}
\toprule
구성요소 & 사례 \\
\midrule
$S$ & 인간 제공 연구 방향(제약된 $S$; 자율적 아이디에이션 없음) \\
$G$ & 인간 아이디어 $\rightarrow$ PhD 에이전트가 문헌 검토로 정교화 \\
$E$ & mle-solver: EDIT/REPLACE + LLM 보상 점수화 + 자기 성찰 + 상위 프로그램 풀 \\
$V$ & $V_\text{semantic}$: NeurIPS LLM 리뷰어; 체계적 과대평가(인간 대비 +2.3점) \\
$M$ & 문헌 컬렉션($M_e$); $M_s$나 $M_p$ 없음 \\
$\pi$ & 코파일럿: 인간; 자율: 고정 순차; 논문 정제 $\rightarrow$ 이전 단계로의 후방 간선 \\
\bottomrule
\end{tabular}
\caption{Agent Laboratory의 \sgmvpi{} 사상.}
\end{table}

\noindent\textbf{(A) 참신성.} 참신성은 인간이 연구 방향을 정함으로써 처리된다. 시스템은 명시적으로 이렇게 진술한다. \textit{``가장 강력한 연구 시스템은 인간이 유도하는 아이디에이션과 LLM 기반 워크플로를 결합할 것이다''}~\citep{schmidgall2025agentlab}.

\noindent\textbf{(B) 루프.} 유일한 후방 간선은 논문 정제에서 이전 단계로 향한다. mle-solver는 분석된 어떤 시스템보다 정교한 $E$ 내부 마이크로 루프를 갖추었다. 종료는 휴리스틱이다(최대 단계 수 또는 타임아웃).

\noindent\textbf{(C) 품질.} 자동 리뷰어는 인간 박사과정 학생 대비 평균 +2.3점만큼 품질을 과대평가한다. 문서화된 환각 사례: GPT-4o가 실제로는 시험된 적 없는 하이퍼파라미터 값(``학습률 0.001, 배치 크기 32, 50 에포크'')을 조작했다.

\subsection{Robin}
\label{sec:robin}

\citet{ghareeb2026robin}은 문헌 탐색 에이전트와 데이터 분석 에이전트를 통합해 과학적 발견의 핵심 지적 단계 전체를 자동화하는 멀티에이전트 시스템을 제안한다. 건성 황반변성(dAMD, dry age-related macular degeneration)이라는 실제 생물의학 과제에 적용해 ripasudil(ROCK 억제제)을 치료 후보로 독자 발굴·검증했다.

\begin{table}[h]
\centering
\small
\begin{tabular}{@{}lp{9cm}@{}}
\toprule
구성요소 & 사례 \\
\midrule
$S$ & dAMD 치료 후보 공간: 기존 약물 재목적화 + 신규 치료 메커니즘(RPE phagocytosis 경로) \\
$G$ & 문헌 탐색 에이전트 $\rightarrow$ 가설 생성 에이전트. RPE phagocytosis 강화 전략 도출 후 ripasudil 제안 \\
$E$ & \textbf{반자율(semi-autonomous) E}: 물리적 세포 실험(RPE 배양·phagocytosis assay)은 인간 연구자 수행; AI는 실험 설계 + RNA-seq 데이터 분석 자율 수행 \\
$V$ & $V_\text{empirical}$: wet-lab 측정값(phagocytosis 개선율, RNA-seq ABCA1 발현 변화). 논문 동료 심사가 $V_\text{human}$ 역할 \\
$M$ & 문헌 데이터베이스 + 실험 결과 아카이브. 각 라운드 결과가 다음 가설 생성에 조건으로 입력 \\
$\pi$ & 반복적 가설-실험-분석 루프. 실험 결과에 따라 다음 치료 전략 방향 갱신 \\
\bottomrule
\end{tabular}
\caption{Robin의 \sgmvpi{} 사상.}
\end{table}

\noindent\textbf{(A) 참신성.} ripasudil의 dAMD 적응증이 기존 문헌에 선례가 없음을 문헌 탐색 에이전트가 확인. 논문 본문의 모든 가설, 실험 계획, 데이터 분석, 그림이 Robin의 산출물이다.

\noindent\textbf{(B) 루프.} 생물의학 도메인 최초의 $V_\text{empirical}$ 분석 루프 폐쇄. 이전 시스템들(Coscientist, ChemCrow)이 분석 측정 단계에서 루프를 열었던 한계를 돌파한다. 물리적 세포 실험(RPE 배양·assay)은 인간 연구자가 수행하며, AI는 실험 설계·RNA-seq 데이터 분석·다음 가설 생성을 자율 수행하는 반자율(semi-autonomous) 구조다. $G \rightarrow V_\text{empirical}(\text{RNA-seq 분석}) \rightarrow G'$의 분석 루프가 AI 주도로 완성된다.

\noindent\textbf{(C) 품질.} 실제 wet-lab 데이터(RPE 세포, phagocytosis assay, RNA-seq)가 $V$의 근거. 현존 AI 과학자 시스템 중 $V_\text{empirical}$이 가장 직접적으로 생물의학 과제에 작동한다. 임상 안전성·유효성 검증은 미래 과제로 남는다.

\subsection{AutoResearchClaw}
\label{sec:autoresearchclaw}

\citet{liu2026autoresearchclaw}은 자기강화 자율 연구 파이프라인 AutoResearchClaw를 제안한다. 다섯 가지 핵심 메커니즘을 도입한다: (1) 가설 정제와 결과 분석을 위한 구조적 다중 에이전트 토론, (2) 실험 실패 시 접근 조정(refine) 또는 전략 전환(pivot)을 선택하는 자기치유 실행기, (3) 수치·인용 전수 검증을 위한 환각 방지 레이어, (4) 과거 실수를 후속 연구 사이클의 안전장치로 변환하는 크로스-프로젝트 실패 메모리, (5) 완전 자율에서 단계별 인간 감독까지의 7단계 인간 개입 모드.

\begin{table}[h]
\centering
\small
\begin{tabular}{@{}lp{9cm}@{}}
\toprule
구성요소 & 사례 \\
\midrule
$S$ & ARC-Bench 25개 토픽의 멀티도메인 연구 아이디어 공간 \\
$G$ & 사전-실험 다중 에이전트 토론으로 가설 정제; 사후 토론으로 결과 분석 \\
$E$ & \textbf{자기치유 실행기}: 실패 시 refine/pivot 결정 루프; 실패를 유용한 데이터로 재활용 \\
$V$ & 환각 검증 레이어(수치·인용 전수 사실확인) + ARC-Bench 외부 벤치마크; wet-lab 없음 \\
$M$ & \textbf{크로스-프로젝트 실패 메모리}: $M_e$ + $M_s$ 통합; 과거 실수를 후속 사이클 guard로 변환 \\
$\pi$ & 7단계 $H$ 모드($H=0$ 완전 자율 $\sim$ $H=1$ 단계별 감독); Goldilocks 구간 실증 \\
\bottomrule
\end{tabular}
\caption{AutoResearchClaw의 \sgmvpi{} 사상.}
\end{table}

\noindent\textbf{(A) 참신성.} 사전-실험 다중 에이전트 토론이 단일 LLM 아이디에이션의 편향을 구조적으로 완화한다. 환각 검증 레이어가 수치·인용을 사실확인해 $V_\text{semantic}$ 신뢰도를 부분 보완한다. 25개 토픽의 ARC-Bench가 외부 비교 기준을 제공한다.

\noindent\textbf{(B) 루프.} 자기치유 실행기와 크로스-프로젝트 메모리가 결합해, 실패가 루프를 끊지 않고 $M_s$로 정제되는 구조를 구현한다. ARC-Bench 25개 토픽에서 AI Scientist v2 대비 54.7\% 향상을 달성했다. 중요한 발견은 인간 개입의 ``골디락스 구간''이다: 완전 자율($H=0$)은 오류를 늘리고, 과도한 감독($H=1$)은 비효율적이며, 고레버리지 결정 지점의 표적 개입이 최선 결과를 산출했다. 이는 \framework{}의 $H \in (0,1)$ 설계 공간에 경험적 증거를 최초로 제공한다.

\noindent\textbf{(C) 품질.} 환각 검증 레이어가 $V_\text{semantic}$ 층의 사실 오류를 구조적으로 차단한다. wet-lab $V_\text{empirical}$은 없으며 벤치마크 점수가 대리 지표다. $H$ 모드 설정에 따라 품질 분산이 존재한다.

\subsection{ScientistOne}
\label{sec:scientistone}

\citet{meng2026scientistone}은 Chain-of-Evidence(CoE) 원리에 기반한 자율 연구 시스템 ScientistOne을 제안한다. 핵심 원칙은 ``모든 주장은 출처에 추적 가능해야 한다''는 것으로, 이를 인용 주장·수치 주장·방법론 주장·결론 주장의 네 범주에 걸쳐 강제한다. 시스템은 세 단계로 구성된다: (1) Semantic Scholar API를 통해 2,000--5,000편의 후보 논문을 수집하고 25--40편의 검증된 참조를 담은 실험 브리프를 생성하는 \textit{Problem Investigator}, (2) $B$ 병렬 분기 $\times$ $I$ 이터레이션의 best-first 탐색으로 솔루션을 탐색하는 \textit{Parallel Explore-Exploit(PEE) Orchestrator}, (3) 주장-증거 인라인 태그가 포함된 연구 표현을 거쳐 \LaTeX{} 원고를 생성하는 \textit{Paper Writer}.

\begin{table}[h]
\centering
\small
\begin{tabular}{@{}lp{9cm}@{}}
\toprule
구성요소 & 사례 \\
\midrule
$S$ & ADRS 벤치마크 5개 시스템 최적화 태스크 + 일반화 6개 태스크; 결정론적 평가기가 암묵적으로 정의 \\
$G$ & Problem Investigator: 2,000--5,000편 $\to$ 500편 $\to$ 25--40편 검증 참조; Ideator가 novelty/feasibility로 순위화된 접근법 제안 \\
$E$ & PEE: $B$ 병렬 분기 $\times$ $E$ 버전/노드 $\times$ $I$ 이터레이션; 스펙 위반 탐지로 분기 프루닝; ablation 자동 실행 \\
$V$ & \textbf{CoE Audit 4중 검증}: I1(점수 재실행)·I2(스펙 위반 LLM 다수결)·I3(참조 API 검증)·I4(코드-방법 정렬); CPR 추가 \\
$M$ & 파일 백업 아티팩트: citation graph($M_s$)·구조화 PDF 노트($M_s$)·실험 로그($M_e$); 작성 단계 연구 표현이 주장-증거 태그로 $M_s$ 연속 유지 \\
$\pi$ & PEE best-first: 상위 $K$ 분기 유지 + 나머지 $B-K$에 신규 아이디에이션; CoE Audit이 $\pi$의 프루닝 신호로 직결 \\
\bottomrule
\end{tabular}
\caption{ScientistOne의 \sgmvpi{} 사상.}
\end{table}

\noindent\textbf{(A) 참신성.} CoE가 G 단계부터 모든 주장의 출처를 강제하므로, 참신성 주장 자체가 기존 문헌과의 검증된 대비에 근거한다. Semantic Scholar 2--5홉 citation graph가 그라운딩을 제공하며 환각 인용 0/337을 달성했다.

\noindent\textbf{(B) 루프.} PEE의 $B \times I$ 탐색 루프가 $V$(CoE Audit)와 직결된다: 스펙 위반 탐지가 분기를 프루닝하고, 최고 점수 분기가 다음 이터레이션 아이디에이션의 조건으로 입력된다. 75편의 논문·5개 시스템·5개 태스크 비교에서 인간 전문가 수준을 달성했다.

\noindent\textbf{(C) 품질.} CoE Audit 4중 검증(I1--I4)으로 환각 인용 0/337, 점수 검증 12/12, 코드-방법 정렬 14/15를 달성했다. 현존 Core System 중 $V_\text{syntax}+V_\text{semantic}$ 조합에서 가장 낮은 환각률. wet-lab $V_\text{empirical}$은 없으며, ADRS 벤치마크가 단일 지표 최적화로 환원된다는 한계가 있다.

\subsection{AutoScientists}
\label{sec:autoscientists}

\citet{gao2026autoscientists}은 분산 자기조직화 멀티에이전트 시스템 AutoScientists를 제안한다. 단일 궤적 또는 중앙집중적 계획 대신 에이전트들이 병렬로 가설을 탐색하고 실험 증거에 따라 적응하며 실패 방향을 공유한다. 세 가지 핵심 메커니즘을 도입한다: (1) 유망 가설 주변으로 에이전트를 자동 재조정하는 \textit{증거 기반 자원 조정}, (2) 자원 배분 전 저품질 가설을 사전 필터링하는 \textit{비평 메커니즘}, (3) 실패 방향을 다른 에이전트의 탐색 조건으로 변환하는 \textit{공유 실패 메모리}. BioML-Bench 24개 생의학 태스크에서 74.4\% 평균 리더보드 백분위(최강 단일 에이전트 대비 +8.33\%), GPT 훈련에서 Autoresearch 대비 1.9배 속도 향상 및 7개 개선 수용(단일 에이전트 0개), ProteinGym 217 assay에서 SOTA 대비 +6.5\%를 보고한다.

\begin{table}[h]
\centering
\small
\begin{tabular}{@{}lp{9cm}@{}}
\toprule
구성요소 & 사례 \\
\midrule
$S$ & 6개 생의학 도메인(영상 세분화·강화·VQA·보고서 생성·병변 검출·단백질 적합도); 공개 벤치마크가 암묵적 정의 \\
$G$ & 에이전트 팀 병렬 가설 제안; 비평 메커니즘이 자원 배분 전 저품질 후보 사전 필터링; 유망 가설 주변 에이전트 자동 재조정 \\
$E$ & 6개 도메인 병렬 계산 실험; 단일 에이전트 대비 1.9$\times$ 속도; wet-lab 없음 \\
$V$ & 공개 리더보드(BioML-Bench·ProteinGym)가 $V_\text{empirical}$; 비평 메커니즘이 사전 $V_\text{semantic}$ 역할 \\
$M$ & 실패 방향 공유($M_e$); 성공 패턴 집단 지식 축적($M_s$ 초기 형태); 에이전트 간 실시간 공유 \\
$\pi$ & 증거 기반 에이전트 재조정: 유망 가설에 자원 집중; 실패 공유로 중복 탐색 회피 \\
\bottomrule
\end{tabular}
\caption{AutoScientists의 \sgmvpi{} 사상.}
\end{table}

\noindent\textbf{(A) 참신성.} 공개 리더보드 성능 우위를 참신성 대리 지표로 활용한다. BioML-Bench·ProteinGym 벤치마크 우위가 발견의 유의성을 대리하며, 참신성 자동 측정은 수행하지 않는다.

\noindent\textbf{(B) 루프.} 분산 자기조직화가 루프 구조다: 병렬 가설 탐색 $\to$ 비평 사전 필터 $\to$ 유망 가설 집중 $\to$ 실험 실행 $\to$ 실패 공유 $\to$ 다음 탐색 갱신. 단일 에이전트 대비 7개 대 0개 개선은 분산 $\pi$가 단일 궤적의 local optima를 회피함을 직접 실증한다.

\noindent\textbf{(C) 품질.} 공개 벤치마크 자동 평가가 $V_\text{empirical}$을 담당한다. wet-lab $V_\text{empirical}$은 없으며 계산 도메인에 한정된다. 공유 실패 메모리($M_e \to M_s$ 변환)는 §\ref{sec:openproblems}의 $M_s$ 미발달 문제에 대한 실증적 부분 응답이나, 자기조직화 창발 행동에 대한 이론적 수렴 보장은 제공되지 않는다.

\subsection{비교 요약}
\label{sec:core-summary}

표~\ref{tab:core-comparison}은 \framework{}이 식별한 핵심 차원을 따라 5개 종단간 시스템을 종합한다.

\begin{table}[h]
\centering
\small
\setlength{\tabcolsep}{3pt}
\begin{tabular}{@{}lcccccc@{}}
\toprule
& \textbf{Coscientist} & \textbf{AI Sci.\ (2024)} & \textbf{AI Sci.\ (2026)} & \textbf{AgentLab} & \textbf{Robin} & \textbf{AutoResearchClaw} \\
\midrule
참신성 척도   & 없음         & LLM 자기 점수 & 3계층         & 인간 위임     & $V_e$(wet-lab) & 다중 에이전트 토론+환각 검증 \\
루프 폐쇄?      & 아니오       & 예(탐욕)      & 예(트리)      & 부분          & 부분($V_e$ 분석) & 예(self-healing) \\
$\pi$ 유형        & ReAct        & 탐욕          & 트리 탐색     & 인간/순차     & 반응적 반복     & 7단계 $H$ 모드 \\
$V$ 계층      & $V_y{+}V_e{+}V_h$ & $V_y{+}V_s^{*}$ & $V_y{+}V_s^{*}{+}V_e$ & $V_y{+}V_s^{*}{+}V_h$ & $V_e$(wet-lab) & $V_s$+환각 검증 \\
$M$ 구조     & 일시적     & $M_e$          & $M_e+M_s$       & $M_e$          & $M_e$            & $M_e+M_s$(크로스-프로젝트) \\
$H$(인간 역할)  & G           & 낮음           & 낮음            & G+$\pi$        & 낮음             & 가변(7단계) \\
도메인      & 화학         & ML             & ML              & ML             & 생물의학(dAMD)  & 멀티도메인 \\
\bottomrule
\end{tabular}
\caption{핵심 AI 과학자 시스템에 대한 \framework{} 비교 분석. $V_y$=구문적, $V_s$=의미적, $V_e$=경험적, $V_h$=인간. 별표는 출처 자신이 계통 편향을 보고한 계층이다. $V$ 계층 행은 존재하는 계층을 나열하며 총체적 완결성 서수가 아니다(코딩 시트에 별도 공개).}
\label{tab:core-comparison}
\end{table}

\subsection{연구 산출물 인프라: Agent-Native Research Artifacts}
\label{sec:ara}

위 4개 시스템은 AI 과학자 루프의 내부 구조를 다룬다. \citet{liu2026lasthuman}은 보완적 문제를 다룬다. 즉, 그 루프의 \emph{출력}이 어떤 형태여야 하며, 어떻게 독립적으로 검증되느냐다. 그들은 서사적 논문을 4계층 파일 시스템 온톨로지로 대체하는 구조적 프로토콜인 Agent-Native Research Artifacts(ARA)를 제안한다. \texttt{/logic}(명시적 반증 기준을 가진 주장, 가설, 관련 연구 의존 그래프), \texttt{/src}(타입 지정 커널 모듈 또는 전체 저장소), \texttt{/trace}(question, decision, experiment, dead\_end, pivot의 5개 타입 노드를 갖는 탐색 그래프), \texttt{/evidence}(원시 기계 가독 결과; 블라인드 재현을 위해 검증 에이전트로부터 보류).

동기는 정량적이다. \citet{liu2026lasthuman}은 논문 중 45.4\%만 재현에 충분한 정보를 담고 있고, 논문 확장 비용의 90.2\%(달러 기준)가 원저자에게 이미 알려진 방법의 실패한 재발견에 소비된다고 보고한다. ARA는 막다른 길과 방향 전환을 포함한 전체 연구 과정을 일급 출판 산출물로 만들어 이 문제를 해결한다.

ARA는 \framework{} 구성요소에 두 가지 직접적 기여를 한다. 첫째, $V$를 ARA Seal로 다룬다. ARA Seal은 $V_\text{syntax}$(스키마 적합성과 계층 간 참조 무결성, 수 초), $V_\text{semantic}$(6차원 루브릭의 Rigor Auditor; 시드 주입된 타당성 변이의 82.6\% 검출), $V_\text{empirical}$(샌드박스 방향성 재현)을 결합한 3단계 자동 검증 자격이다. $V_\text{human}$은 명시적 설계 원칙으로 유의성·참신성·취향을 위해 보존된다: \textit{``유의성, 참신성, 취향에 대한 판단은 인간 리뷰어에게 유보된다''}~\citep{liu2026lasthuman}. 둘째, \texttt{/trace} 계층은 $M_e$와 $M_s$를 다룬다. 포렌식 결합으로 각 주장을 이를 뒷받침하는 실험 및 원시 증거 파일과 연결해, 축적된 지식의 의미 구조를 외재화한다.

핵심적 경험적 발견은 $\pi$에 영향을 미친다. \texttt{/trace} 계층에 보존된 실패 흔적(dead\_ends, pivots)은 확장 과제에서 후속 에이전트의 성능을 가속했지만, 동시에 더 유능한 에이전트를 이전 해결 경로에 고정시켜 제약했다. 더 풍부한 $M_e$가 $\pi$를 단조적으로 개선하지는 않는다.

논문은 ARA를 시스템이 아닌 기반(substrate)으로 위치시킨다: ``ARA는 이런 도구들을 대체하는 것이 아니라 누락된 구조 계층을 제공함으로써 이 간극을 메운다''~\citep{liu2026lasthuman}. \framework{} 용어로 ARA는 $E$의 표준 출력 형식과 $G$의 입력 형식을 명시한다. 그 목표는 AI 과학자 시스템이 서사적 PDF가 아니라 ARA를 생산·소비할 때 축적된 지식이 문장 수준이 아닌 산출물 수준에서 복리로 누적되도록 만드는 데 있다.

\section{도메인 분석}
\label{sec:domains}

우리는 일곱 개 과학 분야에서 \framework{} 구성요소가 어떻게 사례화되는지 검토한다. 각 분야의 $V$ 유형이 그 분야의 검증 병목을 규정하며(명제~\ref{prop:suff}--\ref{prop:novelty}), 게이팅 검증기의 보정 수준은 폐루프 가능성과 공기(共起)한다; 아래 사례들은 이 형식 결과를 예시하는 것이며 열거로 입증하는 것이 아니고, 인과 관계를 확립하지도 않는다.

\subsection{수학}

수학은 독특하게 $V_\text{formal}$을 가능케 한다---Lean 4와 같은 형식 증명 보조기로 구현된 완전한 오라클이다. $\pi$에 깊은 결과가 따른다. 보상 신호가 이진이며 무잡음이므로, 보상 해킹 없이 강화학습이 적용 가능하다. 이 구조는 \citet{vafa2025foundationmodel}이 제기한 더 근본적인 문제, 즉 파운데이션 모델이 훈련 성능에도 불구하고 세계 모델의 귀납적 편향을 내면화하지 못한다는 문제를 수학 도메인이 우회할 수 있는 이유이기도 하다: 형식 커널은 LLM의 세계 모델 이해 여부와 무관하게 수학 공리계를 독립적으로 실행한다. 물리·화학·생물 도메인에서는 이에 상응하는 외부화 메커니즘이 존재하지 않으므로, $G$의 가설 생성뿐 아니라 $E$의 실험 계획과 $V$의 결과 해석 모두 세계 모델 inductive bias의 내면화를 전제한다.

\textbf{FunSearch}~\citep{romera2023funsearch}는 $S$를 Python 함수 공간(우선순위 휴리스틱을 구현하는 짧은 프로그램), $G$를 현재 최선 집단에 조건부인 LLM 변이, $E$를 결정론적 샌드박스 실행, $V_\text{empirical}$을 사용자 제공 평가기의 수치 점수로 사례화한다. 루프는 완전히 닫히며, 참신성은 정량적 SOTA 개선으로 수립된다.

\textbf{DeepSeek-Prover-V1.5}~\citep{xin2024deepseek}는 $S$ = Lean 4 전술 스크립트, $E$ = Lean 증명 커널(결정 가능), $V_\text{formal}$ = 타입 검사 성공(완전 오라클), $\pi$ = RMaxTS(UCB + 내재적 탐색 보상)와 RLPAF(증명 보조기 피드백 기반 강화학습)의 결합을 사용한다. 무잡음 $V$가 보상 해킹 없이 안정적인 강화학습을 가능케 한다.

\textbf{Tsoukalas 등}~\citep{tsoukalas2026formalproof}은 $V_\text{formal}$을 경쟁 수학 벤치마크가 아닌 \emph{실제 미해결 연구 문제}로 처음 확장했다. Erd\H{o}s 문제 목록 353개 중 9개, OEIS 추측 492개 중 44개를 자율 해결했다---비용은 문제당 수백 달러 수준이다. DeepSeek-Prover와의 핵심 차이는 $S$의 개방성에 있다: 이미 형식화된 경쟁 문제가 아닌, 아직 Lean으로 형식화되지 않은 공개 추측을 형식화하고 증명하는 전체 파이프라인을 포함한다. $V_\text{formal}$의 완전 오라클 특성은 인간 개입 없는 완전 폐루프 작동을 가능케 한다.

\textbf{LeanMarathon}~\citep{zhang2026leanmarathon}은 한 걸음 더 나아가 연구 논문 \emph{전체}를 Lean으로 옮기는 장기 autoformalization을 다루며, 병목을 단일 목표 능력이 아닌 \emph{에이전트 내구성}---수 시간 실행 동안 표적을 잃지 않고 일관성을 유지하는 능력---으로 재정의한다. 핵심 추상은 형식 증명 스켈레톤·자연어 증명 그래프·공유 기록을 겸하는 단일 Lean 파일 blueprint이며, 4개 contract-scoped 에이전트가 2단계 오케스트레이터(적대적 표적 충실도 검토 후 동적 증명 DAG의 leaf를 병렬 CI-게이트 방출)로 이를 운영한다. 2026년 두 논문에 걸친 4개 Erd\H{o}s 문제(\#1051, \#1196, \#164, \#1217)에서 7개 표적 정리를 모두 \texttt{sorry} 없이 형식화하고 258개 보조정리·정리를 증명한 반면, 상용 단일 에이전트 기준선은 두 논문 모두에서 실패했다. 프레임워크 관점에서 LeanMarathon은 수학 도메인 내부에서도 $V$가 단일 함수가 아님을 드러낸다: 증명을 입력으로 받던 선행 시스템과 달리 \emph{명제를 형식화}하는 단계에는 형식 오라클이 없으므로, $V$는 $V_\text{formal}$(증명 축, Lean 커널)과 $V_\text{fidelity}$(형식화 충실도 축, Lean 타입이 논문 의도와 일치하는지 감사하는 Target-Reviewer 에이전트)로 분리된다. 즉 autoformalization은 순수 증명이 우회하던 $V_\text{semantic}$ 간극을 표적 명제 계층에 \emph{한정해} 재도입한다.

\textbf{Aletheia}~\citep{google2026aletheia}는 정반대 접근을 취한다. $S$는 자연어 증명 공간, $V$는 오류 가능한 LLM Verifier 서브에이전트, $V_\text{human}$은 진실값으로 남는다. 시도된 700개 Erd\H{o}s 문제 중 인간 감사자는 단 13개(1.9\%)만을 ``의미 있게 정확하다''고 평가했다. 시스템은 명시적으로 ``무의식적 표절''---사전학습 데이터 오염이 발견으로 가장하는 현상---을 핵심적 타당성 위협으로 식별한다.

이 두 패러다임이 정착한 시점에, 범주를 벗어나는 결과가 등장했다. 2026년 5월, OpenAI 내부 범용 추론 모델이 Erd\H{o}s가 1946년에 제기한 평면 단위거리 추측을 반증했다~\citep{openai2026unitdistance}. $n$점 집합의 단위거리 쌍 수가 $n^{1+o(1)}$으로 bound된다는 80년간의 믿음을 깨고, 무한히 많은 $n$에 대해 $\nu(n) \geq n^{1+\delta}$ ($\delta > 0$ 고정)를 만족하는 구성이 존재함을 증명했다. 핵심 아이디어는 대수 수론(무한 비분기 pro-3 tower, Golod--Shafarevich 이론)을 이산 기하 문제에 연결하는 예상치 못한 브리지다. 이 증명은 수학 태스크용으로 특화 훈련하지 않은 모델이 AI가 작성한 문제 진술을 입력 받아 단일 chain-of-thought로 생성했으며~\citep{alon2026unitdistanceremarks}, 외부 수론 전문가들이 정확성을 검증했다. 루프 없음, 반복 없음 — 그럼에도 80년 미해결 문제를 해결했다.

핵심 관찰은 $S$ 선택이 수학 AI 과학자 패러다임을 분기시킨다는 것이다: 프로그램 공간($V_\text{empirical}$로 충분, 완전 자율 가능) vs.\ 자연어 증명 공간($V_\text{human}$ 필요, 부분 자율만 가능). 단위거리 결과는 세 번째 경로를 제시한다: 루프 없는 단일 추론이 전문가 검증을 통과하는 genuine novelty를 생성할 수 있다 --- 적어도 도구적으로 완결된 수학 도메인에서는.

\subsection{물리학}

\textbf{AI Feynman}~\citep{udrescu2020afeynman}은 도메인 지식을 $G$가 아닌 $\pi$에 부호화한다. 탐색 정책은 물리학 특화 휴리스틱(차원 분석, 대칭성 검출, 분리 가능성)을 적용해 $S$의 유효 크기를 점진적으로 축소한다. $M_e$가 부재하다---각 미스터리는 새로 시작한다. 루프는 부분적으로만 닫힌다. 미스터리 간 지식 축적이 없다.

\textbf{Discovering Symbolic Models}~\citep{cranmer2020symbolic}은 2단계 파이프라인을 사용한다. GNN 학습($G$, 상호작용하는 입자 시스템에 대한 귀납 편향) 후 기호 회귀(Eureqa 유전 알고리즘으로 $E$). 루프는 완전히 열려 있다. 발견된 공식은 GNN 학습으로 피드백되지 않는다. $V_\text{human}$이 결정적 역할을 한다: ``이 공식은 우주론자에게 새로운 것이다.''

\textbf{도메인 관찰.} 물리 법칙 발견은 $V$가 수치 정확도와 해석 가능성(기호적 형태) 모두를 평가하기를 요구한다. 이 이중 요구는 알고리즘 발견(정확도만)에는 부재하며, 수학적 구성(SOTA 점수만)보다 더 까다롭다.

\subsection{화학}

\textbf{ChemCrow}~\citep{bran2023chemcrow}은 아키텍처적으로 Coscientist와 유사하다. 18개 API를 갖춘 도구 강화 ReAct이다. 물리적 루프는 분석 측정 경계에서 열린다---로봇 합성기는 반응을 실행하지만 NMR, MS, UV-Vis 특성화는 인간에 위임된다. $V_\text{semantic}$(EvaluatorGPT)은 신뢰할 수 없음이 입증되었다. 명백히 정확한 GPT-4 출력과 명백히 부정확한 출력을 구별하지 못한다.

\textbf{ChemReasoner}~\citep{sprueill2024chemreasoner}은 분석된 어떤 도메인 논문보다 강력한 $M$과 $\pi$를 도입한다. 방문된 모든 (질문, 답변, 보상) 노드를 저장하는 트리 구조 $M_e$, 각 깊이에서 확장 대상으로 최고 점수 노드를 선택하는 빔 탐색 $\pi$(빔 폭 $M=6$)이다. $E$는 계산만 수행하며(GNN 대리자, GemNet-dT), DFT 검증은 사후 검증을 위해 보류된다. 저자들은 GNN 예측과 DFT 진실값 간의 발산을 문서화하여, $V$가 체계적 오차를 보이는 대리자임을 확인한다.

\textbf{도메인 관찰.} 화학의 웻랩 루프는 측정 경계에서 구조적으로 열린다. 물리적 실행은 달성 가능하지만(ChemCrow, Coscientist), 자동화된 분석 특성화는 그렇지 않다. ChemReasoner는 BED에 인접한 $\pi$(보상 순위가 매겨진 노드에 대한 빔 탐색)를 가진 유일한 화학 시스템이다. 최근 Anthropic의 능력 평가~\citep{kamber2026chemist}는 범용 LLM(Claude Opus 4.7)이 NMR 분광 예측에서 전문 소프트웨어(ChemDraw, MestReNova)와 동등하거나 우수한 성능($^1\text{H}$: $\pm0.079$ ppm; 역방향 구조 해명 8/8 완전 성공)을 달성함을 보였다. 이는 루프 개방의 성격을 정밀화한다: 물리적 NMR \emph{측정}은 여전히 장비를 요구하지만, 측정 후 스펙트럼 \emph{해석} 단계는 이제 자동화 가능 영역에 진입했다.

\subsection{재료과학}

재료 발견은 독특한 $V$ 프로파일을 보인다. 밀도 범함수 이론(DFT)으로 계산된 원자당 생성 에너지는 계산적 검증기 가운데 가장 잘 보정된 축에 속한다---계산 비용은 비싸지만 결정적이고 재현 가능하다. 다만 아래에서 보듯 이는 물리적 에너지에 대한 근사이지 완전히 보정된 오라클은 아니다. 한편 물리적 합성은 비용이 많이 들고 인간 매개의 두 번째 검증 계층을 구성한다.

\textbf{GNoME}~\citep{merchant2023gnome}(Graph Networks for Materials Exploration, DeepMind)은 이 구조를 3억 8천만 개의 가설 무기 결정 구조로 확장한다. $G$는 열역학적 안정성을 예측하도록 학습된 두 개의 그래프 신경망(GNN) 쌍이다. 알려진 결정 원형에 작동하는 구조 GNN과 원소 조합에 작동하는 조성 GNN이다. $E$는 고처리량 DFT다. $\pi$는 엔트로피 가중 불확실성 표집이다. GNN이 가장 불확실해하는 후보가 DFT 평가에 우선순위가 부여되며, 이것이 $b(S)$에 대한 능동학습 정책을 근사한다. $M_s$는 GNN이 축적된 DFT 결과로 재학습되면서 성장하며, GNoME은 $G(M_s, S)$가 축적된 실험과 함께 진정으로 개선되는 몇 안 되는 시스템 중 하나다.

\begin{table}[h]
\centering
\small
\begin{tabular}{@{}lp{9cm}@{}}
\toprule
구성요소 & 사례 \\
\midrule
$S$ & 3억 8천만 개의 가설 무기 결정 구조 \\
$G$ & 이중 GNN(구조 + 조성)이 안정성 예측; DFT 결과로 재학습 \\
$E$ & 고처리량 DFT(계산); 협력 연구실에서의 합성(물리) \\
$V$ & $V_\text{empirical}$: 생성 에너지 / 컨벡스 헐 임계(DFT); $V_\text{human}$: 합성 확인 \\
$M$ & $M_e$: 460K 사전 DFT 결과; $M_s$: 매 반복 갱신되는 GNN 가중치 \\
$\pi$ & 불확실성 표집: 고엔트로피 예측을 DFT에 우선순위 부여 \\
\bottomrule
\end{tabular}
\caption{GNoME의 \sgmvpi{} 사상.}
\end{table}

\noindent 핵심 결과: 220만 개의 안정한 구조 식별(83\% 참신, 정밀도 0.87), 그중 736개가 로봇 연구실과의 협업으로 실험적으로 합성되었다. 계산 루프는 닫혀 있다. 합성 루프는 열려 있고, 물리적 검증에는 수 주에서 수 개월의 지연이 따른다. 다만 이 실험적 실현 주장은 비평의 대상이 되어 왔다. \citet{cheetham2024gnome}은 예측된 화합물을 기능성이 실증된 검증된 재료와 동일시해서는 안 된다고 지적한다---이 관점에서 DFT는 완전히 보정된 경험적 오라클이라기보다 물리적 에너지의 근사로 읽어야 하며, 재료 발견의 물리적 루프는 부분적으로만 닫혀 있다.

\textbf{도메인 관찰.} 재료과학은 $\pi$가 BED 이론에 가장 가깝게 근사하는 영역이다. GNoME의 불확실성 표집 정책은 확률적 안정성 모델에 대한 정보 이득 최대화와 구조적으로 동등하며---다른 도메인의 어떤 시스템보다 EIG 최적에 가깝다. 일차적 병목은 DFT 벽시계 시간(후보당 수 시간)으로, 루프 빈도를 제한하지만 폐쇄를 막지는 않는다.

\subsection{생물학}

\textbf{BioPlanner}~\citep{odonoghue2023bioplanner}은 프로토콜 생성 품질의 자동 평가(함수 정확도, 인수 정밀도, 순서 정렬)를 가능케 하는 의사함수(pseudofunction) 추상화에 기여한다. 시스템의 일차적 기여는 $G$-$E$-$V$ 루프 폐쇄가 아니라 $V$의 운영화이다. 웻랩 검증은 성공하지만(\textit{E.\ coli} 배양이 수정 없이 실행됨), 결과가 생성기로 피드백되지 않는다.

\textbf{AlphaFold 2}~\citep{jumper2021alphafold}는 \framework{} 사상에서 명시적으로 제외된다. 이는 가설-실험 루프가 아닌 판별 예측기다. $\pi$가 부재하며(재활용은 고정된 3-반복 정제), $G$가 가설 생성을 수행하지 않으며, 추론 동안 실험적 피드백이 발생하지 않는다. AlphaFold 2는 패턴의 사례라기보다 생물학 AI 과학자 개발을 동기화하는 목표 능력을 나타낸다.

\subsection{사회과학}

사회과학은 LLM이 $E$(인간 피험자 실험)를 대체한다는 점에서 독특하다. 이는 구조적 타당성 위협을 만든다. $G$와 $E$가 동일한 기저 모델을 공유할 때, 실험 실행기가 가설 생성기의 암묵적 편향을 확인하기 쉽다.

\textbf{Automated Social Science}~\citep{manning2024automated}는 LLM 시뮬레이션 실험이 경제 이론과 방향성 일치하는 결과를 산출함을 시연한다. 그러나 LLM은 직접 질의될 때 효과 크기를 체계적으로 13.2배 과대평가하며, 시뮬레이션 기반 추정은 이론에 더 가깝다. 시스템의 $\pi$는 외생 SCM 변수에 대한 완전 요인 설계다. 시나리오 간 전환은 자동화되지 않는다($H$가 이 경계에서 개입한다).

\textbf{Generative Agents}~\citep{park2023generative}는 후속 시스템이 사용하는 E-인프라(메모리 스트림 + 성찰 + 계획 아키텍처)를 제공한다. 그 자체로는 AI 과학자 시스템이 아니다---루프는 사회적으로 그럴듯한 행동을 산출할 뿐 과학적 발견을 산출하지 않는다.

\textbf{Turing Experiments}~\citep{aher2023turing}는 메타 논문으로서, LLM이 사회과학을 위한 타당한 $E$를 구성하는지를 측정한다. 핵심 발견: \emph{초정확도 왜곡(hyper-accuracy distortion)}---더 강력한 LLM은 인간 인구 분포를 시뮬레이션하기에 너무 정확한 답변을 산출한다. 이 결과는 사회과학 AI 과학자 시스템에서 $E$ 충실도에 대한 근본적 제약을 식별한다.

\subsection{인지과학}

\textbf{인지과학 자동 발견}~\citep{jagadish2026cogscidiscovery}은 인지과학 연구 사이클 전체의 in-silico 자동화를 제안한다: (1) LLM이 MDP 문법 기반 실험 패러다임을 생성하고, (2) Centaur 재단 모델이 인구통계·정신질환 지표를 조건부로 합성 행동 데이터를 생성하며, (3) GeCCo 파이프라인이 LLM 기반 프로그램 합성으로 계산 인지 모델을 탐색하고, (4) LLM 비평가가 ``흥미로움(interestingness)'' 신호로 발견을 유도한다. 이 설계는 $G$(LLM 실험 설계)·$E$(합성 행동 시뮬레이션)·$V$(모델 적합도 + 흥미로움 점수)를 in silico로 연결하는 완전 폐쇄 루프를 표방한다.

그러나 두 가지 구조적 취약성이 있다. 첫째, $G$와 $E$가 동일한 LLM 기반 모델을 공유할 때 확인 편향의 위험은 사회과학 도메인과 구조적으로 동일하다~\citep{manning2024automated,aher2023turing}. 둘째, 제안 수준의 비전 논문으로서 실증 결과가 없어 Centaur 시뮬레이션이 실제 인간 행동의 신뢰할 만한 대리자인지 검증되지 않았다; 저자들 자신도 이를 핵심 한계로 명시한다. $\pi$는 단일 스칼라가 아닌 다목적 흥미로움(참신성·압축성·일반화·통합)으로 설계하지만, 이 신호가 game 가능하지 않은지는 미해결 문제다.

\subsection{컴퓨터과학}

\textbf{AutoML-Zero}~\citep{real2020automlzero}은 정규화된 진화 알고리즘으로 원시 수학 연산에서 완전한 ML 알고리즘을 진화시킨다. $S$ = 연산 시퀀스(Setup, Predict, Learn의 세 구성요소), $V_\text{empirical}$ = 프록시 과제에 대한 검증 정확도, $\pi$ = 노화 정규화를 갖춘 토너먼트 선택. 루프는 완전히 닫혀 있고 수백만 세대 동안 실행된다.

\textbf{Algorithm Discovery with LLMs}~\citep{sartori2024algodiscovery}는 LLM을 의미론적 변이 연산자로 사용한다. 구문적 비트 플리핑 대신, LLM이 현재 프로그램의 논리를 이해하고 의미 있는 변형을 제안한다. 그 결과 맹목적 진화 탐색에 비해 탐색 공간을 수십 배 줄인다. \citet{surina2025algodiscovery}는 여기에 강화학습을 결합해, 진화 탐색이 산출한 궤적으로 LLM 변이 연산자 자체를 갱신함으로써 $G$가 탐색 경험과 함께 개선되도록 한다.

\textbf{AutoKernel}~\citep{jaber2026autokernel}은 GPU 커널 최적화를 표적으로 삼아, $S$를 CUDA/Triton 구현 공간으로, $V_\text{empirical}$을 하드웨어에서 측정된 실행시간 가속비로 둔다. LLM 에이전트는 구조화된 최적화 플레이북($M_p$)을 읽고, 현재 커널에 대한 코드 편집을 제안하며, 벤치마킹 전 5단계 정확성 게이트를 실행한다. $\pi$는 keep/revert 정책이다. 정확성 검사를 통과하고 $>$1.01배 가속을 달성한 편집은 커밋되며, 그 외에는 \texttt{git reset}으로 롤백된다. 루프는 야간에 무인으로 실행되며(10시간 실행당 300--400회 반복), 반복당 약 20\% 성공률로 2--5배의 커널 가속을 달성한다.

AutoKernel은 절차 메모리 $M_p$가 일차 시스템 구성요소로 가장 명확하게 작동하는 예다. 최적화 플레이북은 전문가 휴리스틱---루프라인 모델 제약, 메모리 합치(coalescing) 패턴, 타일링 전략---을 구조화된 형식으로 부호화하여, 적용 불가한 전략에 대한 LLM 환각 없이 $G$를 안내한다.

\subsection{도메인 간 요약}

표~\ref{tab:domain-comparison}는 도메인 전반의 \framework{} 분석을 요약한다. 이 패턴은 명제~\ref{prop:suff}--\ref{prop:drift}가 분리한 구분을 예시한다: 검토된 도메인 전반에서 게이팅 검증기의 보정 수준은 $\pi$ 능력 및 폐루프 가능성과 공변하지만, 이 공변이 인과적 우선성을 확립하지는 않는다.

\begin{table}[h]
\centering
\footnotesize
\setlength{\tabcolsep}{3pt}
\begin{tabular}{@{}lccccp{2.5cm}@{}}
\toprule
도메인 & $V$ 유형 & $V$ 신뢰성 & 루프 & $\pi$ & 일차적 병목 \\
\midrule
수학(형식) & $V_\text{formal}$ & 완전 오라클 & 완전 폐쇄 & RL/MCTS & 사전학습 오염 \\
수학(NL) & $V_s+V_h$ & 불완전 & 부분 & 순차 정제 & 환각; 표절 \\
물리학 & $V_e+V_h$ & 중간 & 없음/부분 & 결정론적 & 문제 간 $M$ 부재 \\
화학 & $V_e$(대리) & 중간 & 계산만 폐쇄 & 빔 탐색 & 웻랩 루프 개방 \\
재료 & $V_e$(DFT) & 높음 & 계산 폐쇄 & 불확실성 표집 & DFT 벽시계; 합성 지연 \\
생물학 & $V_e+V_h$ & 낮음--중간 & 부분 & 없음 & $V$ 메트릭 불충분 \\
사회과학 & 통계 검정 & 낮음--중간 & 부분 & 요인 설계 & 순환적 타당성 \\
인지과학 & $V_e$(합성)+$V_h$ & 낮음(미검증) & 설계 제안 & 흥미로움 탐색 & $G$-$E$ 공유 LLM 편향 \\
CS/알고리즘 & $V_e$ & 높음 & 완전 폐쇄 & 진화/LLM & 해석성 부재 \\
\bottomrule
\end{tabular}
\caption{일곱 개 도메인에 걸친 \framework{} 도메인 간 분석. $V_s$=의미적, $V_e$=경험적, $V_h$=인간. $V$ 신뢰성은 진실값과의 일치도를 뜻한다.}
\label{tab:domain-comparison}
\end{table}

\section{미해결 과제}
\label{sec:openproblems}

본 분석에서 모든 시스템과 도메인에 걸쳐 지속되는 네 가지 미해결 과제가 드러난다.

\subsection{참신성의 비자동화 가능성}

본 서베이의 모든 시스템은 동일한 근본적 문제에 직면한다. 참신성은 자동화된 수단으로 신뢰성 있게 평가될 수 없다. Coscientist(메커니즘 없음)에서 The AI Scientist Nature(실제 동료 심사를 포함한 3계층 검증)에 이르는 시스템의 궤적은 동일한 결론으로 수렴한다. ARA~\citep{liu2026lasthuman}는 이를 명시적 설계 원칙으로 만들어, 유의성·참신성·취향을 인간 리뷰어에게 유보하면서 심사의 기계적 계층을 자동화한다.

두 가지 구체적 실패 양상이 문서화되어 있다. (1) \emph{LLM 자기 평가 편향}: 참신성과 품질의 체계적 과대평가로 나타난다(Agent Laboratory에서 +2.3점, The AI Scientist Sakana의 환각 결과, 구조적 안전장치에도 불구하고 ARA 산출물 23개 중 17개에서 관찰된 점수 부풀림). 이 편향은 자기평가에만 국한되지 않는다. \citet{kim2026aireviewers}가 45명 전문가와 함께 Nature-family 논문 82편의 AI 리뷰어 비평 2,960건을 실측한 결과, AI 리뷰어는 정확도 면에서 최상위 인간보다 6--10pp 낮으며 논문에 명시된 내용을 부재로 오주장하는 W3 실패 패턴이 관찰됐다---$V_\text{semantic}$ 층 신뢰 불가의 첫 전문가 직접 검증이다. \citet{liu2026automedench}는 의료 영상 AI 연구 에이전트를 5단계 워크플로우(Plan·Setup·Validate·Inference·Submit)로 평가하는 AutoMedBench를 제안해, 검증 단계 실패(37.7\%)가 제출 실패(38.1\%)와 함께 오류의 대부분을 차지하며 Setup $>$ Validate 능력 비대칭을 실측한다---$V_\text{empirical}$ 자동화의 어려움을 의료 도메인에서 독립적으로 확인한다. 이 신뢰 불가는 참신성 채점에만 국한되지도 않는다: 블랙박스 시스템을 역설계하라는---과학적 추론의 전형적 과제---요구를 받은 LLM은 체계적인 질의 접근이 주어져도 이를 신뢰성 있게 수행하지 못한다~\citep{geng2025reliable}. (2) \emph{사전학습 오염}: 수학에서 가장 첨예하다. 시스템이 학습 데이터에 암묵적으로 들어 있던 결과를 ``발견''하기 때문이다(Aletheia가 Erd\H{o}s 문제에 대해 문서화한 바와 같다).

이 진단에 대한 중요한 반례가 수학 도메인에서 등장했다. OpenAI 범용 추론 모델이 Erd\H{o}s 평면 단위거리 추측을 단일 chain-of-thought로 반증했으며~\citep{openai2026unitdistance}, Fields 메달리스트 Gowers를 포함한 외부 전문가 패널이 그 참신성과 정확성을 검증했다~\citep{alon2026unitdistanceremarks}. 이 결과는 $V_\text{formal}$이 가용한 도메인에서, 즉 결과의 올바름을 원칙적으로 전문가가 검증할 수 있는 수학에서, AI가 genuine novelty를 생성할 수 있음을 보여준다. 한편, 라이덴 선언~\citep{leidendeclaration2026}(서명 사이트 \url{https://leidendeclaration.ai} 기준 2026년 7월 현재 3,000명 이상 서명; 선언문 자체에는 서명자 수가 기재되어 있지 않다)은 AI 수학 능력에 대한 상업적 과장을 경고하며, 개별 성과가 체계적 능력으로 해석되는 것을 경계한다---$V_\text{human}$의 커뮤니티적 위치화 기능이 AI 수학 발견의 유의성 판단에도 여전히 필수적임을 제도 수준에서 재확인한다. 같은 선언은 전문가 감독의 불가결성을 다른 각도에서도 뒷받침한다: AI가 산출한 증명은 인간의 검토 없이는 탐지되지 않는 거의 비가시적인 오류를 품을 수 있으며, 따라서 자동 검사가 가장 성숙한 영역에서조차 어떤 자동 계층도 $V_\text{human}$을 온전히 대체하지 못한다~\citep{leidendeclaration2026}. 그러나 이것이 §\ref{sec:openproblems}의 문제를 해소하지는 않는다. 문제는 $V_\text{formal}$이 없는 생물학·사회과학 도메인에서 AI가 제안한 가설이 사전학습 데이터의 재조합인지 진정한 발견인지를 구별할 수 없다는 데 있다.

ARA Seal~\citep{liu2026lasthuman}은 $V$의 하위 참신성 계층을 자동화하기 위한 가장 구체적인 기존 제안이다. 그 Rigor Auditor는 5가지 주입 유형에 걸쳐 시드 주입된 구조적 위반의 82.6\%를 검출하며, 고립된 실험에서는 22\%의 사각지대가 문서화되어 있다. 이 수치는 자동 검증이 참신성 경계 이전에 어디까지 도달하는지를 보정한다. 경계 자체를 다루려면 사전학습 오염에 강건한 평가 프로토콜이 필요하거나, 생성 모델과 평가 모델 간의 명시적 아키텍처 분리가 필요하다---이런 접근은 대체로 미탐구 상태로 남아 있다.

Science of Science의 계량 연구들은 이 구조적 취약성을 경험적으로 확인한다. \citet{park2023disruptive}는 4,500만 편 논문에서 파괴-강화 지수(CD5)가 1945--2010년 사이 90\% 이상 하락했음을 보였다---문헌 corpus 자체가 이미 consolidating 방향으로 수렴해 있다. \citet{wu2019teams}는 소규모 팀이 오래되고 덜 인용된 문헌을 깊이 탐색하며 파괴적 발견을 하는 반면, 대규모 팀(그리고 다중 에이전트 앙상블)은 최신·인기 문헌에 수렴함을 보였다. \citet{evans2008narrowing}은 전자 아카이브가 인용 다양성을 오히려 줄였음을 보였으며---RAG 기반 $G$가 이 메커니즘을 증폭한다. \citet{uzzi2013atypical}은 최고 임팩트 논문이 높은 관습성(85--95th 백분위)에 소수의 비정형 조합을 더한 구조임을 정량화했다; 단순 참신성 최대화는 임팩트가 낮은 셀에 착지한다. 이 네 결과는 RAG·멀티 에이전트·사전학습 corpus를 가진 AI Scientist가 구조적으로 강화적(CD$\approx{-1}$) 산출물을 생성함을 예측하며, $G$의 올바른 목표가 참신성 최대화가 아니라 ``관습적 기반 + 이질적 소수 조합''의 비대칭 분포임을 지정한다.

참신성 자동화의 한계는 기술적 검증 너머의 층위에도 있다. $G$가 자연어 명제가 아닌 \emph{반증 가능한 구성물(falsifiable construction)}---코드, 물질 구조, 수학적 객체, 생물학적 형태---을 생성할 때, 그 구성물이 과학적 주장이 되려면 연구 커뮤니티가 그것을 해당 인식론적 공간에 위치시켜야 한다. Robin의 경우처럼 $G$가 명시적 명제를 생성하든, ERA나 GNoME처럼 구성물을 생성하든, 수학의 반례처럼 대수적 객체를 생성하든, 그 산출물의 \emph{유의성}은 커뮤니티의 현재 관심사와의 관계 속에서만 성립한다. $V_\text{human}$의 역할은 오류 감지보다 이 인식론적 위치화에 있으며, 이 기능은 사전학습 분포를 공유하지 않는 다양한 사전 분포의 커뮤니티~\citep{taniguchi2024cpc}를 요구한다. AI Scientist가 수학·재료과학·소프트웨어 도메인에서 더 완성된 루프를 달성하는 이유는 이들 도메인에서 커뮤니티가 이미 번역 기준을 형식화해두었기 때문이다---Lean 커널, DFT formation energy, 공개 벤치마크가 $V_\text{human}$의 위치화 기능을 부분적으로 대체한다.

\textit{제안하는 반증 시험:} 사전학습 컷오프 날짜로 훈련 집합 포함 여부가 고정된 출판-대-보류 발견 벤치마크에서, 어떤 자동 검증기가 3개 이상의 서로 다른 도메인에 걸쳐 3인 이상의 도메인 전문가 패널과 참신성 판별에서 대등한 성능---전문가 다수결 대비 코헨 $\kappa \geq 0.80$---을 달성한다면, 이 주장은 반증된다. 오염이 통제된 벤치마크와 전문가 기준선을 구성하는 일 자체가 미해결 과제다. 최근의 한 벤치마크는 오염 통제를 위해 2024년 이후의 발견으로 대상을 제한하는 한 걸음을 내디뎠으나~\citep{liu2025researchbench}, 전문가 기준선 간극을 메우기보다 영감 기반 과제 분해를 평가한다.

\subsection{$\pi$의 BED--실제 간극}

\ref{eq:eid}식은 이론적으로 최적인 탐색 정책을 명시한다: EIG를 최대화하는 실험을 선택한다. 이 정책을 근사할 이론적 인프라는 존재하지만~\citep{foster2019variational,foster2021dad,blau2022rl}, 현재 어떤 AI 과학자 시스템도 이를 구현하지 않는다. 대신:

\begin{itemize}
  \item Coscientist: 반응적($\pi$ = 다음 도구 선택)
  \item AI Scientist Sakana: 아카이브 유사도 회피에 의한 아이디어 선별---본고는 이를 탐욕적 $\pi$로 코딩하나, 출처는 선택 목적함수를 정의하지 않는다
  \item AI Scientist Nature: 최선 우선 트리 탐색(GPT-4o의 노드 점수화)
  \item Agent Laboratory: 순차적(인간 유도)
  \item GNoME: 불확실성 표집($\pi \approx$ GNN 사후분포에 대한 EIG 프록시) --- 가장 가까운 근사이지만, 다루기 쉬운 확률 모델을 가진 도메인에 제한됨
\end{itemize}

가용 이론과 배포된 실제 사이의 간극은 두드러지지만, 부분적 예외가 하나 있다. GNoME~\citep{merchant2023gnome}의 불확실성 표집 정책---GNN 안정성 예측이 가장 큰 엔트로피를 보이는 결정 구조를 우선시---은 확률 모델에 대한 정보 이득 최대화와 구조적으로 동등하다. GNN 불확실성을 사후 엔트로피 $H(\theta \mid b)$의 프록시로 다루어, EIG를 명시적으로 계산하지 않고 근사한다. 이 방식은 재료과학에서 작동한다. $V_\text{empirical}$(DFT)이 잘 보정되어 있고 GNN 사후분포가 다루기 쉽기 때문이다. LLM 기반 시스템에서는 어느 조건도 성립하지 않는다. 가설에 대한 ``사후분포''는 언어 모델 가중치에 암묵적이며, EIG는 $p(\theta \mid y)$에 대한 적분이 필요한데 이는 상당한 아키텍처적 변화 없이는 다루기 어렵다. 정보 이득을 효율적으로 추정하도록 학습된 신경망을 사용하는 분할상환된 EIG 근사~\citep{foster2021dad}는 가설 공간의 다루기 쉬운 모델을 유지할 수 있는 도메인에서 실현 가능한 경로가 된다.

LLM 쪽에서도 진입로가 열리기 시작했다. BED-LLM~\citep{choudhury2026bedllm}은 LLM의 예측 분포로부터 사용 가능한 EIG 사후분포를 유도한다. 다만 그 대상은 실험 설계가 아니라 대화형 질의다---따라서 간극은 불가능성이 아니라 신뢰할 만한 $E$ 및 $V$와의 통합 문제다. 한편 최근의 반증 주도 접근들은 가설 선택을 자유로운 생성이 아니라 실행 가능한 시험에 근거지운다. AIGS~\citep{liu2024aigs}는 자동 반증을 견디고 살아남는지로 가설을 선택하고, Popper~\citep{huang2025popper}는 통계적 오류 제어를 갖춘 순차적 반증 실험을 수행한다. 그러나 둘 중 어느 것도 EIG를 계산하거나 최대화하지 않는다: BED 최적 정책을 실현하는 대신 반증 휴리스틱으로 대체할 뿐이다.

\textit{제안하는 반증 시험:} 어떤 시스템이 닫힌 형식의 대리자도 유한 앙상블 대리자도 허용하지 않는 가설 공간---손으로 명세한 저차원 모델이 아니라 자연어 가설이나 프로그램 텍스트 위에 직접 정의된 사후분포---위의 $p(\theta \mid y)$에 대해 EIG를 명시적으로 계산하고 최대화한다면, 이 간극은 반증된다. 이 조건은 GNoME의 심층 앙상블 정책과 같이 다루기 쉬운 모델에 기댄 대리자를 배제한다.

이 간극의 일부는 나아가 개념적이다. EIG 최적 선택은 이미 표상된 가설 공간을 전제하는 반면, 진정으로 새로운 가설이나 표상이 어떻게 \emph{발생}하는지는 현행 BED 정식화의 밖에 있다~\citep{yanai2019night}---명제~\ref{prop:novelty}가 참신성에 대해 확립한 비식별성의 생성 측 대응물이다.

\subsection{의미 메모리 $M_s$의 미발달}

현재 시스템은 주로 $M_e$(에피소드 메모리: 원시 실험 기록 저장)를 구현한다. $M_s$(의미 메모리: 실험 이력에서 추출된 정제된 패턴과 규칙성)는 대체로 부재하다. 그 결과 1,000개의 완료된 실험을 가진 시스템이 10개를 가진 시스템보다 더 지능적으로 탐색하지 못한다. 축적된 경험이 개선된 $G$나 $\pi$로 정제되지 않기 때문이다.

이 부재의 비용은 재현 격차에서 경험적으로 가시화된다. \citet{liu2026lasthuman}은 출판된 결과 확장 비용의 90.2\%(달러 기준)가 실패한 재발견---원저자가 가졌지만 출판하지 않은 중간 지식의 재구성---에 소비된다고 보고한다. 이것이 정확히 $M_s$가 부호화해야 하지만 그렇지 못한 지식이다: 전문가를 초보자와 구분하는 막다른 길, 방향 전환, 어렵게 얻은 휴리스틱.

ARA \texttt{/trace} 계층~\citep{liu2026lasthuman}은 가장 가까운 기존 근사다. 타입 노드(experiment, dead\_end, pivot)와 증거에 대한 포렌식 결합이 있는 탐색 그래프를 부호화한다. 그러나 이는 학습 메커니즘이 아닌 표현 프레임워크다---$M_s$는 구조화되어 있지만 실험 결과로부터 자동으로 정제되지는 않는다. 경험적으로, 보존된 흔적은 확장 과제에서 에이전트 성능을 가속했지만 더 유능한 에이전트를 이전 해결 경로에 고정시키기도 했다. $M_s$ 구조가 저수준 이력에 대해 선택적으로 망각하거나 추상화하는 메커니즘과 짝지어져야 함을 시사하는 결과다.

외부 $M_s$ 인프라 접근도 병행 시도되고 있다. \citet{qiao2026sciatlas}의 SciAtlas는 4,300만 편 논문·26개 분야·1억 5,700만 엔티티·30억 트리플릿 규모의 다학제 지식 그래프를 구축해, tri-path 협력 검색(키워드·벡터·그래프 경로 병렬 활용)과 그래프 리랭킹으로 단순 RAG를 넘어 위상 추론을 가능하게 한다. 이 접근은 $M_s$를 개별 시스템 내부에 학습하는 대신 공유 외부 인프라로 제공한다는 점에서 \citet{taniguchi2024cpc}의 집합적 공유 표현 $w$와 구조적으로 유사하다. 그러나 이 외부화된 $M_s$가 실제로 $G$의 탐색 품질을 개선하는지에 대한 폐쇄 루프 증거는 아직 확보되지 않았다.

$G(M_s, S)$가 축적된 경험과 함께 진정으로 개선되는 시스템---단순히 아카이브에 추가하는 것이 아니라 전이 가능한 지식을 정제하는 시스템---을 구축하는 일은 장기 과학적 발견에 중요한 함의를 가진 미해결 과제다.

이 병목의 구조적 원인은 탐색 공간에 대한 중간 표현 학습의 부재로 귀착된다. 에이전틱 진화 연구는 이 문제를 다른 각도에서 드러낸다: AlphaEvolve~\citep{georgiev2025alphaevolve}(해 진화), Darwin G\"{o}del Machine~\citep{zhang2025dgm}(에이전트 진화), AEvo~\citep{aevo2026}(절차 진화)는 추상화 수준이 높아질수록 탐색 공간이 커지지만, 세 시스템 모두 LLM 사전학습 표현을 암묵적 매니폴드로 사용하며 탐색 중 표현을 갱신하지 않는다. $M_s$가 진정으로 발달하려면 에피소드 이력으로부터 절차 공간의 압축된 표현을 탐색과 병렬로 학습하는 메커니즘이 필요하다.


\textit{제안하는 반증 시험:} 어떤 시스템이 $n$개의 완료된 실험 이력을 보유한 조건과 이력을
제거한 조건에서 \emph{동일한 기반 모델과 동일한 연산 예산으로} 평가받고, 이전에 보지 않은
과제 집합에서 이력 보유 조건이 유의하게 높은 성공률을 보이거나 유의하게 적은 실험 수로 동등한
성공률에 도달하며, 그 이득이 $n$에 대해 단조 증가한다면, $M_s$ 미발달 진단은 반증된다.
결정적 대조군이 하나 더 필요하다: 원시 이력 전체를 문맥에 넣은 조건이다. 이득이 그 조건에서
재현된다면 관측된 것은 정제된 $M_s$가 아니라 $M_e$ 조회이므로 반증으로 치지 않는다. 이
진단이 주장하는 바는 축적된 경험이 \emph{전이 가능한 형태로 압축}된다는 것이지 검색 가능한
형태로 보존된다는 것이 아니다.
\subsection{실험 충실도 간극}
\label{sec:op4}

시험 시점 연산(test-time compute)은 구조적 비대칭을 만들었다. 가설 생성은 이제 빠르고 저렴해진 반면, 실험 실행 $E$는 여전히 물리적으로 속도가 제한된다---추론을 더 한다고 DFT 벽시계 시간, wet-lab 회전 시간, 합성 지연이 줄지 않는다. 그러나 조사된 어떤 시스템도 계산적 대리 $E$에서 물리 실험으로 \emph{언제} 격상할지에 대한 원리적 기준을 갖고 있지 않다. 결과 캐싱은 과거 결과를 축적하지만 더 높은 충실도의 측정에 언제 투자할지는 결정하지 못한다. 격상 기준이 없으면 계산적 루프 폐쇄는 물리적 발견을 함의하지 않는다.

이 패턴은 도메인 전반에서 일관된다. 재료과학에서 GNoME은 DFT로 38만 개의 안정 화합물을 예측하지만, ICSD에 이미 존재하거나 추가된 구조와 일치하는 것은 736개이며 그중 184개만이 프로젝트 시작 이후의 신규 발견이다~\citep{merchant2023gnome}. 이들은 시스템이 촉발한 합성이 아니라 다른 연구자들이 독립적·동시적으로 만든 구조와의 대조 결과다: 격상 결정이 외부에서 이루어지며 원리적 자동화가 없다. 화학에서 ChemCrow와 ChemReasoner는 계산 루프를 닫지만 wet-lab 루프는 열어 둔다---NMR과 MS 특성 분석은 물리적 실행을 요구하는데, 어떤 현행 시스템도 이를 언제 호출해야 하는지 알지 못한다~\citep{gromski2019chemical}. 생물학에서 BioPlanner의 $V$ 메트릭은 루프를 자율적으로 닫기에 불충분하며~\citep{odonoghue2023bioplanner}, 물리적 실행이라는 물음에는 도달조차 하지 않는다. 인지과학에서 \citet{jagadish2026cogscidiscovery}는 행동 기반 모델을 대리 $E$로 사용하지만(인간 피험자 대신 합성 참가자 데이터), 생체 내 검증으로 격상하는 기준은 제시하지 않으며 스스로 합성 데이터를 근거 진실이 아니라 가설 가속기로 규정한다. GNoME이 가장 성공적인 사례인 것은 정확히 $E$를 단일 충실도(DFT)에 \emph{고정}하고 그 제약 안에서 최적화하기 때문이다: 격상 문제를 푸는 것이 아니라 우회한다. ERA 역시 고정된 계산 평가기 안에서 최적화하여, 물리적 격상 없이 유보된 품질 지표로 판정되는 전문가 수준의 경험적 소프트웨어를 생성한다~\citep{aygun2026era}.

이 문제는 앞의 두 미해결 과제와 구조적으로 구별된다: 완결된 $V$는 자신이 받은 충실도에서만 참신성을 평가하고, EIG 최적 $\pi$는 현재 가용한 충실도의 실험들 가운데서 선택한다. 따라서 $E$ 격상 기준이 없으면 시스템은 신뢰할 만한 과학적 주장이 요구하는 물리적 증거에 이르지 못한 채 계산 루프를 무한히 닫을 수 있다.

최근의 자율 ``자율주행'' 실험실은 이 간극과 직접 관련되지만 이를 메우지는 못한다. 재료과학에서 A-Lab~\citep{szymanski2023alab}은 계산 스크리닝을 로봇 합성에 결합해 처음에는 58개 표적 화합물 중 41개의 합성을 보고했고, 이후 출판 후 회절 재분석을 거쳐 보고된 성공 40건 중 36건 확인(4건 판정 불가)으로 수정됐다~\citep{szymanski2026correction}. 그러나 A-Lab은 물리적 검증이 언제 정당화되는지에 대한 증거 기반 기준이 아니라 고정된 스크리닝-후-합성 규칙으로 격상한다. ORGANA~\citep{darvish2024organa} 같은 로봇 화학 플랫폼은 wet-lab 실험의 물리적 실행을 자동화하고, AutoLabs~\citep{panapitiya2025autolabs} 같은 자기 교정 다중 에이전트 시스템은 하드웨어 실행 가능한 프로토콜의 생성을 자동화하며, 리뷰들은 기반 하드웨어의 빠른 성숙을 기록한다~\citep{tobias2025selfdriving}. 그러나 $E$를 자동화하는 것은 역량을 공급할 뿐 기준을 공급하지 않는다: 이 시스템들은 고정되거나 외부에서 선택된 충실도에서 작동하며, 축적된 계산적 증거가 언제 물리적 검정으로의 격상을 정당화하는지를 스스로 결정하지 않는다. 따라서 이 병목은 하드웨어 자동화와 직교한다---물리 루프가 기계적으로 실행 가능해져도 존속한다.

\textit{제안하는 반증 시험:} 어떤 시스템이 $E$를 계산적 대리에서 물리 실험으로 격상하는 증거 기반 기준을 구현하고, 사전 등록된 벤치마크에서 평가했을 때 고정 격상 시스템보다 물리적 실험 비용 단위당 실험적으로 검증된 발견의 비율이 더 높다면, 이 간극은 반증된다.

\section{논의}
\label{sec:discussion}

\textbf{시스템 설계에 대한 함의.} 본 분석에서 설계 위계가 드러난다. $V$의 선택이 다른 모든 설계 결정에 선행해야 한다. 폐루프 운용이 가능한지와 $\pi$가 어떤 형태를 취할지를 결정하기 때문이다. $V$가 잡음 있는 대리자로 근사되는 도메인(화학, 생물학)은 보상 잡음에 강건한 $\pi$ 전략을 요구한다---빔 탐색이 부분적으로 다루는 요구사항이지만(ChemReasoner) 아직 체계적으로 연구되지 않았다.

\textbf{$H$는 해법인가, 아니면 문제의 핵심인가.} 본 분석은 $H > 0$---어떤 형태로든의 인간 개입---을 현재 AI 과학자 시스템의 설계 필요조건으로 확인한다. 그러나 이를 해법으로 선언하기 전에 더 근본적인 질문을 제기해야 한다. Nature 편집부는 Robin·Co-Scientist·ERA 세 논문과 같은 날 발표한 사설에서 이렇게 묻는다: ``AI가 생산한 지식은 과학적 지식인가?''~\citep{nature2026editorial}

이 질문의 난이도는 $H$ 설정으로 해소되지 않는다. 첫째, \emph{지식 계보의 문제}다. 현재 AI 과학자 시스템의 $G$(가설 생성기)는 수십억 개의 인간 저작 문서로 훈련된 LLM을 사용한다. $G$가 제안한 가설이 진정으로 새로운 것인지, 아니면 훈련 데이터에 산재한 단서들의 비명시적 재조합인지를 판별하는 메커니즘이 없다. Aletheia~\citep{google2026aletheia}는 이를 ``무의식적 표절''로 명명한다. 수학에서 형식 증명 커널이 이를 부분적으로 감지하지만, 생물의학과 사회과학에서는 해당 메커니즘 자체가 부재하다. \citet{vafa2025foundationmodel}은 이 우려에 실증적 기반을 제공한다: 귀납적 편향 프로브(inductive bias probe) 실험에서 궤도 예측으로 훈련된 파운데이션 모델이 뉴턴역학을 새 태스크에 적용하지 못했다---훈련 성능의 우수성이 세계 모델 귀납적 편향의 내면화를 보장하지 않음을 의미한다. 이 문제는 $G$에 국한되지 않는다. 세계 모델 없는 패턴 매칭은 $E$(실험 계획 시 이론적 제약 무시)와 $V$(선행연구·이론과의 인과적 정합성 판단 대신 텍스트 유사도로 대리) 단계에도 동일하게 적용된다---루프 전체가 제1원리가 아닌 훈련 분포 위에 올라서 있다는 구조적 취약성이다.

둘째, \emph{순환 오염의 문제}다. AI 과학자가 생산한 논문이 문헌에 편입되면, 차세대 LLM은 AI가 생성한 내용을 인간 지식으로 학습한다. Zhao et al.~\citep{zhao2026hallucinations}이 실증한 환각 인용의 급증은 이 순환의 초기 신호다. $H > 0$이 이 순환을 막는다고 볼 수 없다. 인간의 검토가 사전학습에서 물려받은 편향을 제거하지는 않기 때문이다. \citet{koo2026kci}가 보인 한국어 학술 생태계의 어휘 변화는 이 침투가 언어 공동체별로 불균등하게 진행됨을 시사한다.

셋째, \emph{정당성의 문제}다. 도구주의적 관점에서는 ripasudil이 dAMD 환자를 실제로 치료한다면 발견의 인식론적 계보가 중요하지 않다고 주장할 수 있다. 그러나 과학은 결과의 집합이 아니라 방법론적 신뢰성에 대한 공동체적 합의에서 권위를 얻는다. $H$를 높게 설정해 인간이 AI의 가설을 선택하고 실험을 승인할 때, 그 지식의 인식론적 권위가 인간에게서 비롯된다고 말할 수 있는가? 이 질문에 대한 답은 우리 프레임워크의 수학적 구조 내에 없으며, 본 논문은 이 경계를 명시적으로 인정한다.

\textbf{집단 AI 공동 과학자를 향하여.} 본 서베이에서 분석한 모든 시스템은 단일 에이전트 루프로 작동한다. 자연스러운 확장은 \emph{집단} AI 과학자다: 공유 외부 의미 메모리 $w$를 공유하고 집단 검증기 $V_\text{collective}$로 조율하는 다중 에이전트 시스템 $\{$\framework$_k\}_{k=1}^K$. \citet{taniguchi2024cpc}는 이 비전을 집단 예측 코딩으로 형식화하며, 다양한 사전 분포 $\theta_k$를 가진 에이전트 커뮤니티가 실험과 메트로폴리스-헤이스팅스 합의 과정으로 모델링된 동료 심사로 공유 표현 $w$를 갱신하는 분산 베이즈 추론으로 과학을 재정식화한다. 그들의 프레임워크에서 사회적 객관성은 사전 분포 다양성에서 나온다---현재 LLM 기반 시스템은 모두 동일한 사전학습 분포를 공유하므로 이 요건을 충족하지 못한다. \citet{taniguchi2024cpc}는 이를 AI의 집단 과학 참여를 가로막는 핵심 장애물인 ``소통 장벽(communication barrier)''으로 규정하며, 이는 우리가 제시한 세 가지를 넘어서는 네 번째 미해결 과제를 구성한다.

특히 우리 분석의 네 가지 미해결 과제는 집단 AI 공동 과학자 실현의 필요조건이기도 하다: $V_\text{collective}$ 구축은 참신성 자동화 해결을 요구하고, 원리적 집단 $\pi$는 BED--실천 간극 해소를 요구하며, 공유 $w$는 발달된 $M_s$를 요구하고, 희소한 물리 검증 역량을 집단이 어떻게 배분할지는 $E$ 격상 기준을 요구한다. 이런 수렴 구조는 네 문제 중 어느 하나에서의 진전이 집단 시스템을 향해 불균형적으로 큰 이득을 가져옴을 시사한다. 구조적 연구 산출물 프로토콜 $w$로서의 ARA~\citep{liu2026lasthuman}는 가장 실현 가능한 진입점이다. 이 비전의 제도적 실험은 이미 시작됐다: Agents4Science(Stanford, 2025)는 AI가 논문 작성과 동료 심사를 모두 수행한 최초의 컨퍼런스로 \emph{Nature}에 보도됐으며~\citep{gibney2025agents4science}, CAISc 2026은 검증 가능한 문제(자동화된 oracle 검증)와 개방형 문제(AI 리뷰 후 인간 검토)의 이중 트랙으로 이를 계승한다. 두 행사는 §\ref{sec:openproblems}의 LLM-as-judge bias 실증 결과와 직접 긴장 관계에 있다---AI 리뷰어 커뮤니티가 사전학습 분포를 공유한 상태에서 \citet{taniguchi2024cpc}가 요구하는 사전 분포 다양성을 달성할 수 있는지가 핵심 미해결 문제로 남는다. 이런 집단은 또한 인식적 다양성을 보존해야 한다. AI 리뷰어는 이미 독립적인 인간 심사자보다 훨씬 많이 겹치므로~\citep{kim2026aireviewers}, 인식 공동체에 이로운 것으로 알려진 일시적 다양성(transient diversity)을 위협한다~\citep{zollman2010epistemic}.

\textbf{에이전틱 진화와 표현 학습.} 에이전틱 진화 연구는 AlphaEvolve~\citep{georgiev2025alphaevolve}(해 직접 진화), Darwin G\"{o}del Machine~\citep{zhang2025dgm}(에이전트 자기수정), AEvo~\citep{aevo2026}(진화 절차의 메타-편집)로 이어지는 계보를 형성하며, 탐색의 추상화 수준을 단계적으로 높여왔다. 이 계보의 이론적 기반은 \citet{risi2025neuroevolution}에 포괄적으로 정리되어 있으며, 특히 목표 함수 없이 다양성을 우선하는 Novelty Search와 Quality-Diversity 알고리즘이 경사하강 기반 최적화가 달성하기 어려운 개방형(open-ended) 창의적 탐색을 가능하게 하는 메커니즘을 다룬다. 이 방향의 공학적 기반으로, \citet{sarkar2025eggroll}은 진화 전략을 수십억 파라미터 모델 스케일로 확장해 GRPO와 경쟁적임을 실증했다. 그러나 이 계보 전체에 공통된 구조적 한계가 있다: 세 시스템 모두 LLM 사전학습 표현을 암묵적 매니폴드로 사용하며, 탐색 공간이 커질수록 증가하는 희소성 문제를 탐색 중 표현 갱신으로 해소하지 않는다. 이는 $M_s$ 미발달의 에이전틱 진화판이다---에피소드 증거($M_e$)는 축적되지만 절차 공간에 대한 전이 가능한 추상 지식으로 정제되지 않는다. 이 병목을 해소하려면 world model 기반 탐색이나 계층적 추상화처럼 탐색과 병렬로 표현을 학습하는 메커니즘이 필요하며, 이는 AI 과학자 시스템의 $\pi$ 개선과도 직결된다.

\textbf{$G$의 출력 스펙트럼.} 프레임워크는 $G$가 자연어 가설을 생성한다고 암묵적으로 가정하지만, 이는 지나치게 좁다. Robin은 명제적 가설을 생성하지만, ERA는 소프트웨어 구성을, GNoME는 결정 구조를, OpenAI 단위거리 반증은 대수적 객체를 생성한다. 이들은 모두 반증 가능한 구성물---수학의 반례와 같은 인식론적 구조---이며, 구성물의 거동이 곧 주장의 검증이다. $G$의 보다 일반적인 정의는 자연어 명제와 반증 가능한 구성물 모두를 포괄한다.

\textbf{한계.} 시스템 간·도메인 간 코딩은 저자에 의한 1차 출처의 정성적 독해로서, 형식 결과를 예시하기 위한 것이지 체계적 경험 연구를 구성하지 않는다; 따라서 우리는 인과 추정이 아니라 연관을 보고하며, 우리의 명제는 명시된 가정(가정~\ref{asm:tract}) 하에서만 성립하고 개방형·비정상 영역은 향후 과제로 남긴다. 우리의 프레임워크 \sgmvpi{}는 또한 중요한 도메인 특화 미묘함을 추상화한다. 도메인에 따라 수십 배로 변하는 $E$의 비용 구조, $S$의 비정상성(개방형 문제), 단일 과제 발견과 다중 과제 발견의 구분이 그것이다. 단일 에이전트 프레임은 또한 \citet{taniguchi2024cpc}가 형식화한 사회적 계층을 추상화한다; \framework{}의 다중 에이전트 설정으로의 확장은 향후 연구의 자연스러운 방향이다.

\section{결론}
\label{sec:conclusion}

우리는 AI 과학자 시스템을 특성화하기 위한 형식적 프레임워크 \framework{}, \sgmvpi{}를 도입하고, 이를 4개의 핵심 종단간 시스템, 1개의 연구 산출물 인프라 제안, 일곱 개 과학 분야의 도메인별 사례에 적용했다. 이 프레임워크는 AI 과학자 문헌을 POMDP와 베이즈 실험 설계의 확립된 이론과 연결하여, 탐색 정책을 평가·개선하는 원리적 기반을 마련한다.

시스템 간, 도메인 간 분석에서 네 가지 관찰이 두드러진다. 첫째, 우리 표본에서 아키텍처적 진보의 대부분---Coscientist의 인간 매개 검증에서 The AI Scientist의 자동 리뷰어를 거쳐 GNoME의 DFT 오라클에 이르기까지---은 $V$ 완결성의 향상과 \emph{함께} 나타난다. 다만 이 시스템들은 $V$를 $G$·$E$·$M$·$\pi$로부터 분리하지 않으므로, 우리는 이를 인과적 우선성이 아니라 반복적으로 관찰되는 병목으로 읽는다. 명제~\ref{prop:suff}--\ref{prop:novelty}는 관측, 보상 검증, 참신성 판단을 왜 분리해야 하는지에 대한 형식적 이유를 제시할 뿐, 완결된 검증기가 폐쇄의 일차적 결정 요인임을 확립하지는 않는다. 서베이는 이 패턴을 열거로 입증하기보다 예시한다. 둘째, $\pi$의 BED--실제 간극은 실재하지만 절대적이지 않다. GNoME의 불확실성 표집 정책은 다루기 쉬운 확률적 환경에서 EIG를 근사하며, 가설 공간의 학습 가능한 모델이 존재하는 도메인에서 이 간극이 해결 가능한 공학적 문제임을 보인다. 셋째, $M_s$의 미발달이 가장 적게 다루어진 병목이다. 시스템은 $M_e$만으로도 루프를 닫지만, 정제 없이는 축적된 지식을 복리로 누적하지 못한다---이것이 AI 과학자 시스템이 비싼 일회성 도구로 남을지 진정으로 누적적인 연구 동반자가 될지를 결정하는 능력이다. 넷째, 위 셋에서 진전이 있더라도 $E$ 격상 기준이 없으면 그 성과는 계산 영역에 머문다.

\noindent\textbf{한계.} \framework{}은 여러 요인을 추상해 버린다: $E$의 비용 구조(도메인 간 수십 배 차이), $S$의 비정상성, 단일 에이전트 프레이밍이 놓치는 사회적 계층~\citep{taniguchi2024cpc}. 더 근본적으로, $V$ 완결성 패턴은 확립된 인과적 결정 요인이 아니라 교차 시스템 \emph{연관}이며, 검토된 시스템들은 선택 편향과 생존 편향에 노출된 편의 표집이다---실패해 논문화되지 않은 시스템은 애초에 관측되지 않는다. 게다가 기계적 루프 폐쇄는 완결성을 요구하지 않으므로(AI Scientist v1과 NovelSeek은 낮은 완결성에서 계산 루프를 닫는다), 이 주장은 \emph{신뢰할 만한} 폐쇄에 관한 것인데 우리는 이를 $V$와 독립적으로 측정하지 않았다. 검증기가 신뢰할 수 없다고 우리가 주장하는 시스템들을 LLM 코더로 특성화한 것은 반사적 한계이며, 그래서 중심 주장을 보정 하위 구성개념에 한정하고 인간 전문가 코딩을 향후 과제로 남긴다.

우리가 식별한 네 가지 미해결 과제---참신성 비자동화 가능성, BED--실제 간극, $M_s$ 미발달, 실험 충실도 간극---는 독립적이지 않다. 완전한 $V$와 다루기 쉬운 $M_s$를 가진 시스템은 구성에 의해 EIG 최적 $\pi$를 구현할 수 있을 것이다. 그리고 앞의 세 문제가 모두 해결되더라도 $E$ 격상 기준이 없으면 그 성과는 계산 영역 안에 머문다: $V$는 자신이 받은 충실도에서만 판정하고 $\pi$는 현재 가용한 충실도에서만 선택하므로, 실험 충실도 간극은 나머지 셋의 진전을 물리적 발견으로 옮기는 마지막 관문이다. 네 문제 모두 동일한 구조적 공백으로 수렴한다는 점은 이 분야에 단일 조직 질문을 던진다. \textit{어떻게 폐루프 운용을 지원하기에 충분히 완전하면서 동시에 원리적 탐색을 지원하기에 충분히 구조화된 검증기와 메모리 구조를 구축하는가?}

\bibliographystyle{abbrvnat}
\bibliography{references_full}

\end{document}
